% Producing slides that display illustrative and attractive presentations — ICT Upper Sixth — Cours signé OnBuch+
% Official MINESEC syllabus. Compiler avec : tectonic ICT22-producing-slides-that-display-illustrative-and-attractive-pr.tex
\newcommand{\DOCMATIERE}{ICT}
\newcommand{\DOCNIVEAU}{Upper Sixth}
\newcommand{\DOCMODULE}{Upper Sixth — ICT}
\newcommand{\DOCLECON}{22}
\newcommand{\DOCTITRE}{Producing slides that display illustrative and attractive presentations}
% OnBuch+ — préambule « cours signés » ENRICHI (compilé avec tectonic / XeLaTeX)
% Système visuel « Pop » d'OnBuch+ : fond crème (jamais blanc), bordures encre
% épaisses, ombres « dures » décalées, titres Archivo Black en capitales.
% Version MATHÉMATIQUES : graphes (pgfplots/tikz), boîtes pédagogiques
% (définition, propriété, méthode, attention, le savais-tu, exercice résolu,
% exercice, TP, bilan), page de garde et signature OnBuch+ ancrée en bas.
%
% Macros attendues dans main.tex AVANT \input{preamble} :
%   \DOCMATIERE  \DOCNIVEAU  \DOCMODULE  \DOCLECON (numéro)  \DOCTITRE
\documentclass[11pt]{article}

\usepackage{fontspec}
\usepackage[english]{babel}
\usepackage[a4paper,margin=2cm,top=3.4cm,bottom=2.8cm,headheight=1.4cm,footskip=1.3cm]{geometry}
\usepackage{amsmath,amssymb,amsfonts}
\usepackage{mathtools} % \xmapsto, \coloneqq... souvent utilisés par les modèles
\usepackage{cancel} % simplifications barrées (bilans de forces)
% \abs{z} (module, valeur absolue) et \norm{u} (norme), utilisés par les modèles
\providecommand{\abs}{}\let\abs\relax\DeclarePairedDelimiter{\abs}{\lvert}{\rvert}
\providecommand{\norm}{}\let\norm\relax\DeclarePairedDelimiter{\norm}{\lVert}{\rVert}
\usepackage[table]{xcolor} % [table] : fournit aussi \rowcolors (alternance de lignes)
\usepackage{pagecolor}
\usepackage{tikz}
\usetikzlibrary{shapes.symbols,circuits.logic.US,circuits.logic.IEC,arrows.meta,positioning,calc,shapes.geometric,shapes.misc,shapes.arrows,decorations.pathmorphing,decorations.pathreplacing,decorations.markings,patterns,shadows,mindmap,trees,fit,backgrounds,matrix,intersections,angles,quotes}
\usepackage{pgfplots}
\usepackage[version=4]{mhchem} % équations nucléaires \ce{^{235}_{92}U}
\usepackage[european,siunitx]{circuitikz} % schémas électriques
\pgfplotsset{compat=1.18}
\usepgfplotslibrary{fillbetween,groupplots,polar,statistics,dateplot} % aires entre courbes (calcul intégral)
\usepackage{enumitem}
\usepackage{fancyhdr}
% [most] charge toutes les bibliothèques tcolorbox (listings, minted, raster,
% documentation...) et gonfle inutilement la mémoire TeX sur un document long
% (~300 boîtes) : on ne charge que ce qui est réellement utilisé.
\usepackage[skins,breakable]{tcolorbox}
\usepackage{graphicx}
\usepackage{colortbl}
\usepackage{booktabs}
\usepackage{tabularx}
\usepackage{diagbox} % case d'en-tête diagonale des tableaux à double entrée
% Colonnes extensibles centrée / gauche / droite (C, L, R), souvent utilisées par les modèles
\newcolumntype{C}{>{\centering\arraybackslash}X}
\newcolumntype{L}{>{\raggedright\arraybackslash}X}
\newcolumntype{R}{>{\raggedleft\arraybackslash}X}
\usepackage{array}
\usepackage{multicol}
\usepackage{siunitx}
% parse-numbers=false : accepte aussi \SI{2,5 \times 10^{-3}}{...}
\sisetup{locale=UK,output-decimal-marker={.},group-minimum-digits=4,parse-numbers=false}
\DeclareSIUnit\ohm{\ensuremath{\symup{\Omega}}} % U+2126 absent de la police
\DeclareSIUnit\division{div} % réglages d'oscilloscope (ms/div, V/div)
\DeclareSIUnit\year{yr}\DeclareSIUnit\annum{yr}\DeclareSIUnit\Ma{Ma}\DeclareSIUnit\tr{tr} % tours (vitesse de rotation, ex. \SI{1500}{\tr\per\minute})
\usepackage{qrcode}
% \degree : commande fréquemment utilisée par les modèles pour un angle ou une
% température en dehors de \SI{}{} (non fournie par les paquets chargés).
\providecommand{\degree}{^{\circ}}
% \qed (fin de démonstration), utilisé par les modèles sans amsthm
\providecommand{\qed}{\hfill$\square$}
% \barre : double barre des valeurs interdites dans un tableau de variations (tkz-tab)
\providecommand{\barre}{\ensuremath{\|}}
% environnement proof (amsthm non chargé) utilisé par les modèles
\ifdefined\proof\else\newenvironment{proof}[1][Proof]{\par\noindent\textit{#1.}\ }{\qed\par}\fi
\usepackage{lastpage}
\usepackage{hyperref}
\hypersetup{hidelinks}

% ============================================================
% Palette « Pop » OnBuch+ — mêmes hex que la classe P (plus_ui.dart)
% ============================================================
\definecolor{poporange}{HTML}{F59321}
\definecolor{popdark}{HTML}{E07A0C}
\definecolor{popcream}{HTML}{FFF6E8}
\definecolor{popink}{HTML}{1C1714}
\definecolor{poppurple}{HTML}{7A5AE0}
\definecolor{popgreen}{HTML}{2E9E6A}
\definecolor{poppink}{HTML}{E0457B}
\definecolor{popblue}{HTML}{2F7DE1}
\definecolor{popgold}{HTML}{C98A12}
\definecolor{popmuted}{HTML}{8A7F76}
\definecolor{popline}{HTML}{EADFD2}
\definecolor{popgray}{HTML}{8A7F76}
\colorlet{popgrey}{popgray}
% Alias tolérants (le modèle nomme parfois les couleurs différemment)
\colorlet{popwhite}{white}
\colorlet{popblack}{popink}
\colorlet{popred}{poppink}
\colorlet{popyellow}{popgold}
% Teintes claires pour fonds de tableaux / zones
\colorlet{popblueL}{popblue!10!white}
\colorlet{poporangeL}{poporange!14!white}
\colorlet{popgreenL}{popgreen!12!white}
\colorlet{poppurpleL}{poppurple!12!white}
\colorlet{poppinkL}{poppink!12!white}
\colorlet{popgoldL}{popgold!14!white}

\pagecolor{popcream}

% ============================================================
% Typographie
% ============================================================
\defaultfontfeatures{Ligatures=TeX}
\setmainfont{PlusJakartaSans-Medium.ttf}[Path=../../../fonts/,BoldFont=PlusJakartaSans-Bold.ttf]
% Maths en sans-serif (Fira Math) avec chiffres et lettres droites de Plus
% Jakarta, pour que les formules se fondent dans le texte.
\usepackage[math-style=ISO,bold-style=ISO]{unicode-math}
\setmathfont{FiraMath-Regular.otf}[Path=../../../fonts/]
\setmathfont{PlusJakartaSans-Medium.ttf}[Path=../../../fonts/,range={up/num,up/{Latin,latin},\mathup}]
% (icomma retiré : décimales avec point en anglais)
% Caractères mathématiques Unicode absents des polices de texte (Plus Jakarta,
% Archivo Black) : rendus par leur équivalent mathématique, en texte comme en
% formule — sinon ils s'affichent en carré vide (ex. « Arithmétique dans ℤ »).
\usepackage{newunicodechar}
\newunicodechar{ℝ}{\ensuremath{\mathbb{R}}}
\newunicodechar{ℤ}{\ensuremath{\mathbb{Z}}}
\newunicodechar{ℂ}{\ensuremath{\mathbb{C}}}
\newunicodechar{ℕ}{\ensuremath{\mathbb{N}}}
\newunicodechar{ℚ}{\ensuremath{\mathbb{Q}}}
\newunicodechar{𝔻}{\ensuremath{\mathbb{D}}}
\newunicodechar{α}{\ensuremath{\alpha}}
\newunicodechar{β}{\ensuremath{\beta}}
\newunicodechar{γ}{\ensuremath{\gamma}}
\newunicodechar{δ}{\ensuremath{\delta}}
\newunicodechar{ε}{\ensuremath{\varepsilon}}
\newunicodechar{θ}{\ensuremath{\theta}}
\newunicodechar{λ}{\ensuremath{\lambda}}
\newunicodechar{μ}{\ensuremath{\mu}}
\newunicodechar{σ}{\ensuremath{\sigma}}
\newunicodechar{τ}{\ensuremath{\tau}}
\newunicodechar{φ}{\ensuremath{\varphi}}
\newunicodechar{ψ}{\ensuremath{\psi}}
\newunicodechar{ω}{\ensuremath{\omega}}
\newunicodechar{Σ}{\ensuremath{\Sigma}}
% majuscules produites par \MakeUppercase dans les titres (α-aminés -> Α)
\newunicodechar{Α}{\ensuremath{\alpha}}
\newunicodechar{Β}{\ensuremath{\beta}}
\newunicodechar{Γ}{\ensuremath{\Gamma}}
\newunicodechar{Δ}{\ensuremath{\Delta}}
\newunicodechar{Ε}{\ensuremath{\varepsilon}}
\newunicodechar{Θ}{\ensuremath{\Theta}}
\newunicodechar{Λ}{\ensuremath{\Lambda}}
\newunicodechar{Μ}{\ensuremath{\mu}}
\newunicodechar{Τ}{\ensuremath{\tau}}
\newunicodechar{Φ}{\ensuremath{\Phi}}
\newunicodechar{Ψ}{\ensuremath{\Psi}}
\newunicodechar{Ω}{\ensuremath{\Omega}}
\newunicodechar{↦}{\ensuremath{\mapsto}}
\newunicodechar{∈}{\ensuremath{\in}}
\newunicodechar{∉}{\ensuremath{\notin}}
\newunicodechar{∧}{\ensuremath{\wedge}}
\newunicodechar{⇔}{\ensuremath{\Leftrightarrow}}
\newunicodechar{⟺}{\ensuremath{\Longleftrightarrow}}
\newunicodechar{⇒}{\ensuremath{\Rightarrow}}
\newunicodechar{⟹}{\ensuremath{\Longrightarrow}}
\newunicodechar{∘}{\ensuremath{\circ}}
\newunicodechar{∪}{\ensuremath{\cup}}
\newunicodechar{∩}{\ensuremath{\cap}}
\newunicodechar{≡}{\ensuremath{\equiv}}
\newunicodechar{⊂}{\ensuremath{\subset}}
\newunicodechar{⊥}{\ensuremath{\perp}}
\newunicodechar{∃}{\ensuremath{\exists}}
\newunicodechar{∀}{\ensuremath{\forall}}
\newunicodechar{∛}{\ensuremath{\sqrt[3]{\,}}}
\newunicodechar{∜}{\ensuremath{\sqrt[4]{\,}}}
\newunicodechar{⃗}{\ensuremath{{}^{\to}}}
\newunicodechar{ⁿ}{\ifmmode^{n}\else\textsuperscript{\ensuremath{n}}\fi}
\newunicodechar{ˣ}{\ifmmode^{x}\else\textsuperscript{\ensuremath{x}}\fi}
\newunicodechar{ᵇ}{\ifmmode^{b}\else\textsuperscript{\ensuremath{b}}\fi}
\newunicodechar{ʳ}{\ifmmode^{r}\else\textsuperscript{\ensuremath{r}}\fi}
\newunicodechar{ᵘ}{\ifmmode^{u}\else\textsuperscript{\ensuremath{u}}\fi}
\newunicodechar{ʸ}{\ifmmode^{y}\else\textsuperscript{\ensuremath{y}}\fi}
\newunicodechar{ᵃ}{\ifmmode^{a}\else\textsuperscript{\ensuremath{a}}\fi}
\newunicodechar{ᵉ}{\ifmmode^{e}\else\textsuperscript{\ensuremath{e}}\fi}
\newunicodechar{ᵏ}{\ifmmode^{k}\else\textsuperscript{\ensuremath{k}}\fi}
\newunicodechar{ᵖ}{\ifmmode^{p}\else\textsuperscript{\ensuremath{p}}\fi}
\newunicodechar{ᴺ}{\ifmmode^{N}\else\textsuperscript{\ensuremath{N}}\fi}
\newunicodechar{ᶜ}{\ifmmode^{c}\else\textsuperscript{\ensuremath{c}}\fi}
\newunicodechar{ʰ}{\ifmmode^{h}\else\textsuperscript{\ensuremath{h}}\fi}
\newunicodechar{ᵠ}{\ifmmode^{q}\else\textsuperscript{\ensuremath{q}}\fi}
\newunicodechar{ₙ}{\ifmmode_{n}\else\textsubscript{\ensuremath{n}}\fi}
\newunicodechar{ₐ}{\ifmmode_{a}\else\textsubscript{\ensuremath{a}}\fi}
\newunicodechar{ᵢ}{\ifmmode_{i}\else\textsubscript{\ensuremath{i}}\fi}
\newunicodechar{ᵣ}{\ifmmode_{r}\else\textsubscript{\ensuremath{r}}\fi}
\newunicodechar{ₖ}{\ifmmode_{k}\else\textsubscript{\ensuremath{k}}\fi}
\newunicodechar{ᵦ}{\ifmmode_{\beta}\else\textsubscript{\ensuremath{\beta}}\fi}
\newunicodechar{ₓ}{\ifmmode_{x}\else\textsubscript{\ensuremath{x}}\fi}
\newfontfamily\popdisplay{ArchivoBlack-Regular.ttf}[Path=../../../fonts/]
\newfontfamily\poplogo{SpaceGrotesk-Bold.ttf}[Path=../../../fonts/]
\newfontfamily\popsemi{PlusJakartaSans-SemiBold.ttf}[Path=../../../fonts/]
\newfontfamily\popxbold{PlusJakartaSans-ExtraBold.ttf}[Path=../../../fonts/]
\newfontfamily\popxboldls{PlusJakartaSans-ExtraBold.ttf}[Path=../../../fonts/,LetterSpace=3.0]

\setlist[enumerate]{leftmargin=1.7em,itemsep=5pt,topsep=4pt}
\setlist[itemize]{leftmargin=1.7em,itemsep=4pt,topsep=4pt,label={\color{poporange}\textbullet}}
\setlength{\parindent}{0pt}
\setlength{\parskip}{5pt}
\renewcommand{\arraystretch}{1.25}

% pgfplots : style Pop par défaut
\pgfplotsset{
  every axis/.append style={axis line style={popink,line width=1pt},tick style={popink},
    grid style={popline,line width=0.5pt},label style={font=\small},tick label style={font=\footnotesize}},
  popcurve/.style={poporange,line width=1.8pt,no markers,smooth},
}
% Même style disponible dans un \draw TikZ simple (hors pgfplots)
\tikzset{popcurve/.style={poporange,line width=1.8pt}}
\tikzset{popgrid/.style={popline,very thin}}
\colorlet{popcurve}{poporange} % le nom du style est parfois utilisé comme couleur

% ============================================================
% Pill / squiggle
% ============================================================
\newcommand{\poppill}[3]{%
  \tikz[baseline=-0.55ex]\node[draw=popink,line width=1.2pt,rounded corners=2.6pt,%
    fill=#2,inner xsep=6.5pt,inner ysep=3pt]{%
    \popxboldls\fontsize{7}{7}\selectfont\color{#3}\MakeUppercase{#1}};%
}
\newcommand{\popsquiggle}[1]{%
  \tikz[baseline=-0.4ex]{
    \draw[#1,line width=2.6pt,line cap=round]
      plot [smooth] coordinates {(0,0.09) (0.34,0.01) (0.68,0.21) (1.02,0.01) (1.36,0.10)};
  }%
}
% Mot-clé surligné (marqueur orange sous le texte)
\newcommand{\cle}[1]{{\popxbold\color{popdark}#1}}

% ============================================================
% En-tête / pied
% ============================================================
\pagestyle{fancy}
\fancyhf{}
\renewcommand{\headrulewidth}{0pt}
\renewcommand{\footrulewidth}{0pt}
\lhead{\raisebox{-0.15ex}{\poplogo\fontsize{13}{13}\selectfont\color{popink}OnBuch\color{poporange}+}}
\rhead{\poppill{\DOCMATIERE\ \textbullet\ \DOCNIVEAU\ \textbullet\ Chapter \DOCLECON}{white}{popink}}
\lfoot{\popxbold\fontsize{7.6}{7.6}\selectfont\color{popmuted}\MakeUppercase{Signed course OnBuch+}\ \textbullet\ {\popsemi\DOCTITRE}}
\rfoot{\tikz[baseline=-0.6ex]\node[rounded corners=8.5pt,fill=popink,minimum height=17pt,minimum width=17pt,inner xsep=5pt,inner sep=0]{\popxbold\fontsize{8}{8}\selectfont\color{popcream}\thepage\,/\,\pageref*{LastPage}};}

% ============================================================
% Titres
% ============================================================
\newcommand{\chaptitre}[2]{%
  \begin{tcolorbox}[enhanced,colback=poporange,colframe=popink,boxrule=2.6pt,arc=11pt,
      drop shadow=popink,left=15pt,right=15pt,top=12pt,bottom=11pt,before skip=3pt,after skip=18pt]
    \poppill{#2}{popink}{popcream}\par\vspace{9pt}
    {\raggedright\hyphenpenalty=10000\exhyphenpenalty=10000\popdisplay\fontsize{22}{24}\selectfont\color{popcream}\MakeUppercase{#1}\par}
  \end{tcolorbox}
}
% Section numérotée : pastille orange avec le numéro + titre Archivo
\newcounter{popsec}
\newcommand{\coursec}[1]{%
  \stepcounter{popsec}\setcounter{popsub}{0}%
  \par\vspace{16pt}\noindent
  \begin{minipage}{\linewidth}
  \tikz[baseline=-0.7ex]\node[draw=popink,line width=1.6pt,fill=poporange,rounded corners=4pt,minimum size=22pt,inner sep=0]{\popdisplay\fontsize{12}{12}\selectfont\color{popcream}\thepopsec};%
  \hspace{8pt}{\popdisplay\fontsize{15}{17}\selectfont\color{popink}\MakeUppercase{#1}}\par
  \vspace{4pt}\noindent{\color{poporange}\rule{36pt}{3.6pt}}
  \end{minipage}\par\nopagebreak\vspace{8pt}
}
\newcounter{popsub}
\newcommand{\courssub}[1]{%
  \stepcounter{popsub}%
  \par\vspace{9pt}\noindent{\popxbold\fontsize{11.5}{13}\selectfont\color{popink}{\color{poporange}\thepopsec.\thepopsub}\hspace{6pt}#1}\par\nopagebreak\vspace{4pt}
}

% ============================================================
% Boîtes pédagogiques — même recette : fond blanc, bordure épaisse colorée,
% ombre dure encre, badge plein. Titre optionnel [#1].
% ============================================================
\tcbset{popbase/.style={breakable,enhanced,colback=white,boxrule=2.3pt,arc=9pt,
  shadow={3pt}{-3pt}{0pt}{popink},left=14pt,right=14pt,top=11pt,bottom=11pt,before skip=11pt,after skip=15pt,
  fontupper=\popsemi\fontsize{10.6}{14.5}\selectfont\color{popink},
  fontlower=\popsemi\fontsize{10.6}{14.5}\selectfont\color{popink}}}
% Coche dessinée (le glyphe ✓ n'existe pas dans les polices du document)
\newcommand{\popcheck}[1]{\tikz[baseline=-0.5ex]\draw[#1,line width=1.6pt,line cap=round,line join=round](-0.13,0.0)--(-0.04,-0.09)--(0.13,0.1);}
\providecommand{\coche}{\popcheck{popgreen}}
\providecommand{\resultat}[1]{\par\smallskip\noindent\fcolorbox{popgreen}{popgreenL}{\parbox{\dimexpr\linewidth-2\fboxsep-2\fboxrule\relax}{\textbf{Result.} #1}}\par\smallskip}
\newcommand{\popboxhead}[3]{\poppill{#1}{#2}{white}\ifx\relax#3\relax\else\hspace{6pt}{\popxbold\fontsize{10.6}{12}\selectfont\color{popink}#3}\fi\par\vspace{7pt}}

\newtcolorbox{aretenir}[1][]{popbase,colframe=popgreen,before upper={\popboxhead{Key points}{popgreen}{#1}}}
\newtcolorbox{exemplebox}[1][]{popbase,colframe=poporange,before upper={\popboxhead{Exemple}{poporange}{#1}}}
\newtcolorbox{definition}[1][]{popbase,colframe=popblue,before upper={\popboxhead{Definition}{popblue}{#1}}}
\newtcolorbox{propriete}[1][]{popbase,colframe=poppurple,before upper={\popboxhead{Property}{poppurple}{#1}}}
\newtcolorbox{methode}[1][]{popbase,colframe=popgold,before upper={\popboxhead{Method}{popgold}{#1}}}
\newtcolorbox{attention}[1][]{popbase,colframe=poppink,before upper={\popboxhead{Warning}{poppink}{#1}}}
\newtcolorbox{experience}[1][]{popbase,colframe=popblue,colback=popblueL,before upper={\popboxhead{Experiment}{popblue}{#1}}}
% « Le savais-tu ? » : carte violette pleine (contexte, culture, Cameroun)
\newtcolorbox{savaistu}[1][]{popbase,colback=poppurple,colframe=popink,
  fontupper=\popsemi\fontsize{10.4}{14.2}\selectfont\color{white},
  before upper={\poppill{Did you know?}{white}{poppurple}\ifx\relax#1\relax\else\hspace{6pt}{\popxbold\color{white}#1}\fi\par\vspace{7pt}}}
% Exercice résolu : énoncé en haut, \tcblower puis la solution sur fond crème
\newcounter{popexo}\newcounter{popexr}
\newtcolorbox{exoresolu}[1][]{popbase,colframe=poporange,colbacklower=popcream,
  segmentation style={popink,line width=1.2pt,dashed},
  before upper={\stepcounter{popexr}\popboxhead{Worked example \thepopexr}{poporange}{#1}},
  before lower={\poppill{Solution}{popink}{popcream}\par\vspace{6pt}}}
% Exercice (non résolu en place) — la correction est dans la partie Corrigés
\newtcolorbox{exercice}[1][]{popbase,colframe=popink,
  before upper={\stepcounter{popexo}\popboxhead{Exercise \thepopexo}{popink}{#1}}}
% Corrigé
\newtcolorbox{corrige}[1][]{popbase,colframe=popgreen,colback=popgreenL,
  before upper={\popboxhead{Solution}{popgreen}{#1}}}
% Objectifs / prérequis / bilan
\newtcolorbox{objectifs}{popbase,colframe=popink,colback=poporangeL,
  before upper={\popboxhead{Chapter objectives}{popink}{}}}
\newtcolorbox{prerequis}{popbase,colframe=popink,colback=popblueL,
  before upper={\popboxhead{Prerequisites}{popblue}{}}}
\newtcolorbox{bilan}{popbase,colframe=popink,colback=white,boxrule=2.8pt,
  before upper={\popboxhead{Fiche bilan}{poporange}{L'essentiel à connaître pour l'examen}}}
% Figure encadrée avec légende
\newtcolorbox{popfigure}[1][]{enhanced,breakable=false,colback=white,colframe=popline,boxrule=1.4pt,arc=8pt,
  left=10pt,right=10pt,top=10pt,bottom=8pt,before skip=10pt,after skip=12pt,halign=center,
  lower separated=false,halign lower=center,
  fontlower=\popsemi\fontsize{9}{11.5}\selectfont\color{popmuted},
  #1}
\newcommand{\legende}[1]{\par\vspace{6pt}{\popsemi\fontsize{9}{11.5}\selectfont\color{popmuted}#1}}

% ============================================================
% Signature OnBuch+
% ============================================================
\newcommand{\poplogoinline}[1]{{\poplogo\fontsize{#1}{#1}\selectfont\color{popink}OnBuch\color{poporange}+}}
\newcommand{\signatureonbuch}{%
  \begin{tcolorbox}[enhanced,colback=popink,colframe=popink,boxrule=2.2pt,arc=10pt,
      drop shadow=poporange,left=14pt,right=14pt,top=10pt,bottom=10pt,before skip=0pt,after skip=0pt]
    \tikz[baseline=-0.7ex]\node[circle,fill=popgreen,draw=popcream,line width=1.3pt,minimum size=22pt,inner sep=0]{\popcheck{white}};%
    \hspace{9pt}\begin{minipage}[c]{0.62\linewidth}
      {\popxbold\fontsize{12}{13}\selectfont\color{popcream}Signed\ \textbullet\ {\poplogo OnBuch\color{poporange}+}}\par\vspace{2pt}
      {\raggedright\popsemi\fontsize{8}{10.5}\selectfont\color{popline}\MakeUppercase{OnBuch+ teaching team}\\\MakeUppercase{Follows the official MINESEC syllabus}\par}
    \end{minipage}\hfill
    {\poplogo\fontsize{22}{22}\selectfont\color{popcream}OnBuch\color{poporange}+}
  \end{tcolorbox}%
}

% ============================================================
% Page de garde
% #1 titre · #2 matière · #3 niveau · #4 module · #5 n° leçon · #6 durée ·
% #7 description / objectifs (texte ou itemize)
% ============================================================
\newcommand{\pagegarde}[7]{%
  \thispagestyle{empty}
  \newgeometry{a4paper,margin=1.7cm,top=1.6cm,bottom=1.4cm}%
  \begin{tikzpicture}[remember picture,overlay]
    \fill[poporange!18] ([xshift=-3.2cm,yshift=-3.6cm]current page.north east) circle (4.6cm);
    \fill[poppurple!14] ([xshift=2.4cm,yshift=7.6cm]current page.south west) circle (3.8cm);
  \end{tikzpicture}%
  {\poplogo\fontsize{30}{30}\selectfont\color{popink}OnBuch\color{poporange}+}\hfill
  \raisebox{0.6ex}{\poppill{Signed course \textbullet\ Premium}{popink}{popcream}}\par
  \vspace{1pt}\popsquiggle{poppurple}\par
  \vspace{8pt}
  \begin{tcolorbox}[enhanced,colback=poporange,colframe=popink,boxrule=2.8pt,arc=13pt,
      drop shadow=popink,left=18pt,right=18pt,top=12pt,bottom=13pt,after skip=10pt]
    \poppill{#2 \textbullet\ #3}{popink}{popcream}\hspace{5pt}\poppill{#4}{white}{popink}\par\vspace{8pt}
    {\popdisplay\fontsize{14}{14}\selectfont\color{popink}CHAPTER #5}\par\vspace{4pt}
    {\raggedright\hyphenpenalty=10000\exhyphenpenalty=10000\popdisplay\fontsize{25}{27}\selectfont\color{popcream}\MakeUppercase{#1}\par}
    \vspace{8pt}
    {\popxbold\fontsize{9.5}{9.5}\selectfont\color{popink}\MakeUppercase{Suggested time: #6}}
  \end{tcolorbox}
  \begin{tcolorbox}[enhanced,colback=white,colframe=popink,boxrule=2.2pt,arc=10pt,
      drop shadow=popink,left=16pt,right=16pt,top=10pt,bottom=10pt,after skip=8pt]
    \poppill{In this chapter}{poppurple}{white}\par\vspace{5pt}
    \popsemi\fontsize{9.3}{12.6}\selectfont\color{popink}#7
  \end{tcolorbox}
  \noindent\begin{minipage}[t]{0.487\linewidth}
    \begin{tcolorbox}[enhanced,colback=popgreen,colframe=popink,boxrule=2.2pt,arc=10pt,
        drop shadow=popink,left=11pt,right=11pt,top=10pt,bottom=10pt,height=3.1cm,valign=center]
      \tcbox[colback=white,colframe=popink,boxrule=1.3pt,arc=5pt,boxsep=0pt,nobeforeafter,
        left=3pt,right=3pt,top=3pt,bottom=3pt,valign=center]{\qrcode[height=1.9cm]{https://play.google.com/store/apps/details?id=com.luvvix.onbuch&hl=en}}%
      \hspace{8pt}\begin{minipage}[c]{0.5\linewidth}
        {\popdisplay\fontsize{11}{12.5}\selectfont\color{white}\MakeUppercase{Download OnBuch}}\par\vspace{4pt}
        {\raggedright\popsemi\fontsize{8.4}{11}\selectfont\color{white}Free \textbullet\ Google Play\\AI tutor, past papers, results\par}
      \end{minipage}
    \end{tcolorbox}
  \end{minipage}\hfill
  \begin{minipage}[t]{0.487\linewidth}
    \begin{tcolorbox}[enhanced,colback=poppurple,colframe=popink,boxrule=2.2pt,arc=10pt,
        drop shadow=popink,left=11pt,right=11pt,top=10pt,bottom=10pt,height=3.1cm,valign=center]
      \tikz[baseline=-0.7ex]\node[circle,fill=white,draw=popink,line width=1.3pt,minimum size=1.6cm,inner sep=0]{\popdisplay\fontsize{24}{24}\selectfont\color{poppurple}+};%
      \hspace{8pt}\begin{minipage}[c]{0.6\linewidth}
        {\popdisplay\fontsize{11}{12.5}\selectfont\color{white}\MakeUppercase{Go OnBuch+}}\par\vspace{4pt}
        {\raggedright\popsemi\fontsize{8.4}{11}\selectfont\color{white}All signed courses, unlimited AI tutor, PDF sheets — OnBuch+ tab in the app\par}
      \end{minipage}
    \end{tcolorbox}
  \end{minipage}
  \vspace{8pt}
  \popcontenu
  \vfill
  \signatureonbuch
  \restoregeometry
  \newpage
}

% Rangée « Ce cours contient » (page de garde)
\newcommand{\poptile}[3]{% #1 couleur #2 gros texte #3 légende
  \begin{tcolorbox}[enhanced,colback=#1,colframe=popink,boxrule=2pt,arc=9pt,shadow={2.5pt}{-2.5pt}{0pt}{popink},
      left=6pt,right=6pt,top=9pt,bottom=9pt,halign=center,valign=center,height=1.75cm,width=\linewidth,nobeforeafter]
    {\popdisplay\fontsize{12.5}{13}\selectfont\color{white}\MakeUppercase{#2}}\par\vspace{3pt}
    {\centering\popsemi\fontsize{7.8}{9.5}\selectfont\color{white}#3\par}
  \end{tcolorbox}}
\newcommand{\popcontenu}{%
  \par\noindent{\popxboldls\fontsize{7.5}{7.5}\selectfont\color{popmuted}THIS SIGNED COURSE CONTAINS}\par\vspace{7pt}
  \noindent\begin{minipage}[t]{0.235\linewidth}\poptile{popblue}{Course}{complete, illustrated, step by step}\end{minipage}\hfill
  \begin{minipage}[t]{0.235\linewidth}\poptile{poporange}{Methods}{and worked examples}\end{minipage}\hfill
  \begin{minipage}[t]{0.235\linewidth}\poptile{poppink}{Exercises}{exam style + solutions}\end{minipage}\hfill
  \begin{minipage}[t]{0.235\linewidth}\poptile{popgreen}{Summary sheet}{the essentials to remember}\end{minipage}\par}

% Sommaire
\newtcolorbox{sommaire}{enhanced,colback=white,colframe=popink,boxrule=2.2pt,arc=10pt,
  drop shadow=popink,left=18pt,right=18pt,top=13pt,bottom=13pt,before skip=6pt,after skip=20pt,
  before upper={\poppill{Contents}{poppurple}{white}\par\vspace{9pt}\popsemi\fontsize{10.6}{14.5}\selectfont\color{popink}}}

% Signature de fin de cours — poussée tout en bas de la dernière page
\newcommand{\findecours}{%
  \par\vfill\nopagebreak
  \begin{center}\popsquiggle{poporange}\end{center}\vspace{4pt}
  \signatureonbuch
}
\AtBeginDocument{\def\llbracket{\mathopen{[\mkern-3.5mu[}}\def\rrbracket{\mathclose{]\mkern-3.5mu]}}} % intervalles d'entiers (glyphe ⟦ absent de la police)
% Glyphes absents de Fira Math : redéfinis à partir de glyphes disponibles
\AtBeginDocument{%
  \def\setminus{\mathbin{\backslash}}\let\smallsetminus\setminus
  \def\vdots{\mathord{\vbox{\baselineskip3.2pt\lineskiplimit0pt\kern4pt\hbox{.}\hbox{.}\hbox{.}}}}%
  \def\ddots{\mathinner{\mkern1mu\raise6.5pt\hbox{.}\mkern2mu\raise3.6pt\hbox{.}\mkern2mu\raise0.7pt\hbox{.}\mkern1mu}}%
  \def\triangle{\mathord{\scalebox{0.9}{$\symup{\Delta}$}}}%
  \def\nabla{\mathord{\raisebox{\depth}{\scalebox{1}[-1]{$\symup{\Delta}$}}}}%
  \def\checkmark{\popcheck{popdark}}%
  \def\star{\mathbin{\ast}}\def\bigstar{\mathord{\ast}}%
  \def\diamond{\mathbin{\scalebox{0.6}{$\Diamond$}}}%
  \def\bigcup{\mathop{\mathchoice{\raisebox{-0.3em}{\scalebox{1.7}{$\cup$}}}{\scalebox{1.25}{$\cup$}}{\cup}{\cup}}\displaylimits}%
  \def\bigcap{\mathop{\mathchoice{\raisebox{-0.3em}{\scalebox{1.7}{$\cap$}}}{\scalebox{1.25}{$\cap$}}{\cap}{\cap}}\displaylimits}%
  \def\lesseqqgtr{\mathrel{\lessgtr}}\def\bigoplus{\mathop{\oplus}\displaylimits}%
}
\providecommand{\ssi}{if and only if} % abréviation fréquente
\AtBeginDocument{\providecommand{\wideparen}{\overparen}} % arc de cercle

% Fonctions usuelles que les modèles utilisent sans que amsmath les fournisse
\providecommand{\arsinh}{\operatorname{arsinh}}\providecommand{\arcosh}{\operatorname{arcosh}}
\providecommand{\artanh}{\operatorname{artanh}}\providecommand{\arsech}{\operatorname{arsech}}
\providecommand{\arcsch}{\operatorname{arcsch}}\providecommand{\arcoth}{\operatorname{arcoth}}
\providecommand{\arcsinh}{\operatorname{arsinh}}\providecommand{\arccosh}{\operatorname{arcosh}}
\providecommand{\arctanh}{\operatorname{artanh}}\providecommand{\sech}{\operatorname{sech}}
\providecommand{\csch}{\operatorname{csch}}\providecommand{\arcsec}{\operatorname{arcsec}}
\providecommand{\arccsc}{\operatorname{arccsc}}\providecommand{\arccot}{\operatorname{arccot}}
\providecommand{\sgn}{\operatorname{sgn}}\providecommand{\lcm}{\operatorname{lcm}}
\providecommand{\Var}{\operatorname{Var}}\providecommand{\Cov}{\operatorname{Cov}}
\providecommand{\permil}{\ensuremath{{}^{0}\!/_{00}}}\providecommand{\textperthousand}{\ensuremath{{}^{0}\!/_{00}}}

%% FIGFIX-BEGIN v5 — correctifs globaux des figures (docs/onbuch-generation/tools/figfix)
\usepackage{adjustbox}\usepackage{etoolbox}\tcbuselibrary{raster}
% toute figure TikZ/pgfplots plus large que la ligne est réduite à la largeur disponible
\BeforeBeginEnvironment{tikzpicture}{\begin{adjustbox}{max width=\linewidth}}
\AfterEndEnvironment{tikzpicture}{\end{adjustbox}}
% légendes pgfplots lisibles par-dessus les courbes
\pgfplotsset{every axis/.append style={legend style={fill=white,fill opacity=0.92,text opacity=1,draw=black!15,inner sep=3pt}}}
% étiquettes d'arêtes (syntaxe edge["..."]) avec fond blanc
\tikzset{every edge quotes/.append style={fill=white,inner sep=1.2pt}}
%% FIGFIX-END
\begin{document}

\pagegarde{Producing slides that display illustrative and attractive presentations}{ICT}{Upper Sixth}{Upper Sixth — ICT}{22}{3 periods}{This chapter teaches Upper Sixth students how to plan, produce and deliver attractive slide presentations using presentation software such as Microsoft PowerPoint, LibreOffice Impress or Google Slides. Learners master slide design, multimedia integration, transitions, animations and rehearsal techniques needed for the GCE Advanced Level practical examination and real-life communication.
\vspace{4pt}
\begin{multicols}{2}\small
\begin{itemize}
\item By the end of this chapter I can explain the purpose and components of presentation software and its interface.
\item By the end of this chapter I can plan a presentation by defining the audience, objective, outline and number of slides.
\item By the end of this chapter I can create, save, duplicate, reorder and delete slides using appropriate layouts.
\item By the end of this chapter I can apply themes, background styles, master slides and consistent formatting to a presentation.
\item By the end of this chapter I can insert and format text, pictures, shapes, tables, charts, audio and video on slides.
\item By the end of this chapter I can apply and configure slide transitions with timing and sound options.
\end{itemize}\end{multicols}}

\begin{sommaire}
\begin{multicols}{2}
\begin{enumerate}
\item Section 1: Getting Started with Presentation Software and Producing Slides
\item Section 2: Designing Attractive and Consistent Slides
\item Section 3: Multimedia, Transitions and Animations
\item Section 4: Delivering, Printing and Evaluating the Presentation
\item Integration activity
\item Methods and worked examples
\item Exercises
\item Solutions to the exercises
\item Summary sheet
\end{enumerate}
\end{multicols}
\end{sommaire}

\begin{objectifs}
\begin{itemize}
\item By the end of this chapter I can explain the purpose and components of presentation software and its interface.
\item By the end of this chapter I can plan a presentation by defining the audience, objective, outline and number of slides.
\item By the end of this chapter I can create, save, duplicate, reorder and delete slides using appropriate layouts.
\item By the end of this chapter I can apply themes, background styles, master slides and consistent formatting to a presentation.
\item By the end of this chapter I can insert and format text, pictures, shapes, tables, charts, audio and video on slides.
\item By the end of this chapter I can apply and configure slide transitions with timing and sound options.
\item By the end of this chapter I can apply entrance, emphasis, exit and motion-path animations to objects and control their order and timing.
\item By the end of this chapter I can run a slideshow, use presenter tools, rehearse timings and print handouts.
\item By the end of this chapter I can evaluate a presentation against design principles (readability, consistency, the 6x6 rule) and improve it.
\item By the end of this chapter I can produce a complete multimedia presentation on a real topic and defend it before an audience.
\end{itemize}
\end{objectifs}

\begin{prerequis}
\begin{itemize}
\item Basic file management: creating folders, saving, opening and renaming files (Lower Sixth ICT).
\item Word processing skills: entering and formatting text, inserting pictures and tables (Lower Sixth ICT).
\item Familiarity with the keyboard, mouse and the general window/menu interface of application software.
\item Notions of charts and data representation from spreadsheet work (Lower Sixth ICT).
\item Basic knowledge of multimedia file types (images, audio, video) seen in the chapter on multimedia.
\end{itemize}
\end{prerequis}

\newpage

\coursec{Section 1: Getting Started with Presentation Software and Producing Slides}

During the end-of-year ceremony at GBHS Bamenda, a student is asked to present the school's results to parents and the Divisional Delegate of Secondary Education. He reads long paragraphs from a sheet of paper, head down, while the audience loses interest after five minutes. His classmate, however, projects colourful slides with photos of the prize winners, animated charts of the pass rates and smooth transitions between topics --- the hall stays attentive to the very end and the Delegate warmly congratulates the school. \textbf{Same information, completely different impact!} In a Cameroon where GCE project defences, business pitches in Douala, church announcements in Yaoundé and political campaign meetings all rely on projected slides, knowing how to produce an \cle{illustrative} and \cle{attractive} presentation is no longer optional. This chapter teaches you to design, enrich and deliver professional slide presentations, and this first section lays the foundations: what presentation software is, how to find your way in its window, and how to create, organise and save your very first slides.

\courssub{Presentation Software: Definition, Purpose and Examples}

\begin{definition}[Presentation software]
\cle{Presentation software} is an application program used to create a sequence of \cle{slides} --- individual screens combining text, images, charts, sound and video --- designed to be projected or displayed to an audience in order to support an oral presentation.
\end{definition}

The key intuition is simple: a slide is \emph{not} a document to be read silently like a Word page; it is a \textbf{visual support} for what the speaker says. The audience listens to you and \emph{glances} at the slide. Good presentation software therefore makes it easy to combine short texts, pictures and multimedia, to arrange them in a logical order, and to project them full-screen.

\begin{center}
\begin{tabularx}{\linewidth}{|l|X|X|}
\hline
\rowcolor{popblueL} \textbf{Software} & \textbf{Publisher / type} & \textbf{Remarks} \\
\hline
Microsoft PowerPoint & Microsoft, commercial (Office 365) & The most widespread in Cameroonian schools and companies; files \texttt{.pptx}. \\
\hline
LibreOffice Impress & The Document Foundation, free and open source & Installed in many school computer rooms; native format \texttt{.odp}. \\
\hline
Google Slides & Google, free, web-based & Runs in a browser; ideal for group work and automatic online saving. \\
\hline
WPS Presentation & Kingsoft, free with premium options & Light, popular on laptops and phones; reads \texttt{.pptx}. \\
\hline
\end{tabularx}
\end{center}

\begin{savaistu}[Presentations in Cameroonian life]
Slide presentations are everywhere: teachers project lessons in multimedia rooms, Upper Sixth students defend their projects before a jury, start-ups in Douala's ``Silicon Mountain'' (Buea) pitch their ideas to investors, churches announce programmes between services, and ministries present budgets with projected charts. Mastering presentation software is a genuine employability skill.
\end{savaistu}

\courssub{Launching the Software and Touring the Interface}

To launch PowerPoint, click the \textbf{Start} button, type \emph{PowerPoint} and press \textbf{Enter} (in LibreOffice, open the Start Centre and choose \emph{Impress}). Whatever the software, the working window is organised around the same zones.

\begin{popfigure}
\begin{tikzpicture}[scale=1.0, every node/.style={font=\small}]
% Title bar
\draw[fill=popblue] (0,7.6) rectangle (12,8.2);
\node[text=white, anchor=west] at (0.15,7.9) {\textbf{Title bar: MyPresentation.pptx --- PowerPoint}};
% Ribbon
\draw[fill=popblueL] (0,6.4) rectangle (12,7.6);
\node[anchor=west, text=popdark] at (0.15,7.35) {\textbf{Ribbon / Tabs:} File \quad Home \quad Insert \quad Design \quad Transitions \quad Slide Show};
\node[anchor=west, text=popmuted] at (0.15,6.75) {Groups of commands: Font, Paragraph, Drawing, Images, Charts \dots};
% Outline pane
\draw[fill=white, draw=popline] (0,1.2) rectangle (2.6,6.4);
\node[align=center, text=popdark] at (1.3,5.9) {\textbf{Outline /}\\ \textbf{Slide pane}};
\node[align=center, text=popmuted, font=\footnotesize] at (1.3,3.6) {Thumbnails\\ of slides\\ 1, 2, 3 \dots};
% Slide pane
\draw[fill=white, draw=popline] (2.6,1.2) rectangle (12,6.4);
\node[text=popdark] at (7.3,5.9) {\textbf{Slide pane (main editing area)}};
\draw[fill=popgreenL, draw=popgreen] (4.2,4.4) rectangle (10.4,5.5);
\node[text=popdark] at (7.3,4.95) {Click to add title};
\draw[fill=white, draw=popgreen] (4.2,2.0) rectangle (10.4,4.1);
\node[text=popmuted] at (7.3,3.05) {Click to add text (placeholder)};
% Notes pane
\draw[fill=popgreenL, draw=popline] (0,0.5) rectangle (12,1.2);
\node[anchor=west, text=popdark] at (0.15,0.85) {\textbf{Notes pane:} speaker notes, visible to you but not projected};
% Status bar
\draw[fill=popmuted!40, draw=popline] (0,0) rectangle (12,0.5);
\node[anchor=west, text=popdark, font=\footnotesize] at (0.15,0.25) {\textbf{Status bar:} Slide 1 of 5 \quad Theme: Ion \quad View buttons \quad Zoom 100\%};
\end{tikzpicture}
\legende{The main zones of a presentation software window (PowerPoint / Impress).}
\end{popfigure}

\begin{aretenir}[Zones of the window]
\begin{itemize}
\item \textbf{Title bar}: name of the open file and of the application.
\item \textbf{Ribbon} (or menus/toolbars in Impress): tabs such as \emph{Home, Insert, Design, Transitions, Slide Show}, grouping all commands.
\item \textbf{Slide pane}: the large central area where you edit the current slide.
\item \textbf{Outline / thumbnail pane} (left): miniature views of all slides, used to select, move or reorganise them.
\item \textbf{Notes pane} (bottom): private speaker notes attached to each slide.
\item \textbf{Status bar and view buttons}: slide number, theme, zoom, and quick access to \emph{Normal}, \emph{Slide Sorter}, \emph{Reading} and \emph{Slide Show} views.
\end{itemize}
\end{aretenir}

\courssub{Creating a New Presentation: Blank or Template}

When the software opens, it offers two starting points:

\begin{itemize}
\item a \textbf{blank presentation}: white slides with no decoration --- total freedom, but you must design everything yourself;
\item a \textbf{template} or \textbf{design theme}: a ready-made set of colours, fonts, background and slide masters created by professionals.
\end{itemize}

\begin{methode}[Creating a presentation with a theme]
\begin{enumerate}
\item Launch the software and choose \emph{New}.
\item Select \emph{Blank presentation} or click a template thumbnail.
\item If you started blank, open the \textbf{Design} tab and click a \cle{theme}: every slide instantly receives the same colours, fonts and background.
\item Type the title and subtitle of the first slide, then save immediately (\textbf{Ctrl+S}).
\end{enumerate}
\end{methode}

\begin{attention}[One theme, not five]
Applying different themes or random colours to each slide makes a presentation look amateurish. Choose \textbf{one} theme for the whole presentation; consistency is the first sign of professionalism.
\end{attention}

\courssub{Slides, Layouts and Placeholders}

A \cle{slide} is one screen of the presentation. A \cle{placeholder} is a dotted box on a slide that ``reserves a place'' for a specific content: a title, a bulleted list, a picture, a chart or a table. You simply click it and type or insert.

\begin{definition}[Slide layout]
A \cle{slide layout} is a predefined arrangement of placeholders (title, text, content) on a slide. Examples: \emph{Title Slide}, \emph{Title and Content}, \emph{Two Content}, \emph{Comparison}, \emph{Blank}.
\end{definition}

\begin{center}
\begin{tabularx}{\linewidth}{|l|X|}
\hline
\rowcolor{popblueL} \textbf{Layout} & \textbf{Typical use} \\
\hline
Title Slide & First slide: presentation title, your name, class, date. \\
\hline
Title and Content & The workhorse: a title plus one list, picture, table or chart. \\
\hline
Two Content & Compare two ideas, or place text beside an image. \\
\hline
Comparison & Side-by-side comparison with titles for each column. \\
\hline
Blank & Full freedom, e.g.\ for a large photo or a diagram. \\
\hline
\end{tabularx}
\end{center}

To change the layout of an existing slide: \textbf{Home} $\rightarrow$ \textbf{Layout} $\rightarrow$ click the desired arrangement. The text already typed is kept and simply rearranged.

\courssub{Managing Slides: Add, Duplicate, Move, Hide, Delete}

All slide management happens in the left thumbnail pane or, better, in the \textbf{Slide Sorter view}, which displays all slides as a grid of miniatures --- perfect for checking the logical flow of your story.

\begin{itemize}
\item \textbf{Add}: \textbf{Ctrl+M} or \emph{Home $\rightarrow$ New Slide}; the new slide appears after the selected one.
\item \textbf{Duplicate}: right-click a thumbnail $\rightarrow$ \emph{Duplicate Slide} (useful to reuse a structure).
\item \textbf{Move}: drag the thumbnail to its new position (in Normal or Slide Sorter view).
\item \textbf{Hide}: right-click $\rightarrow$ \emph{Hide Slide}; the slide stays in the file but is skipped during the show --- handy for backup slides in a defence.
\item \textbf{Delete}: select the thumbnail and press \textbf{Delete}.
\end{itemize}

\begin{exemplebox}[Reorganising a defence presentation]
Manka'a prepares her project defence with slides in this order: Title, Results, Introduction, Method, Conclusion. In Slide Sorter view she notices the illogical order, drags \emph{Introduction} to position 2 and \emph{Method} to position 3, then hides a detailed calculations slide that she will show only if the jury asks. Her 8-slide story now flows naturally.
\end{exemplebox}

\courssub{Entering and Formatting Text}

Click a placeholder and type. Then select the text and use the \textbf{Home} tab to set:

\begin{itemize}
\item \textbf{Font}: choose a clear sans-serif font (Calibri, Arial); avoid fancy scripts that are hard to read from the back of a hall.
\item \textbf{Size}: titles around 36--44 pt, body text at least 24 pt for projection.
\item \textbf{Colour}: strong contrast with the background (dark text on light background or the reverse).
\item \textbf{Alignment} and \textbf{bullets/numbering}: bullets for unordered ideas, numbers for steps or rankings.
\item \textbf{Line spacing}: 1.15 or 1.5 for readability.
\end{itemize}

\begin{attention}[The paragraph trap --- apply the 6x6 rule]
The most common beginner's mistake is typing whole paragraphs on a slide. The audience then reads instead of listening, and the speaker is tempted to read the screen aloud. Apply the \cle{6x6 rule}: about \textbf{6 lines per slide, 6 words per line}. Put only keywords on the slide and \emph{speak} the details --- that is exactly what the notes pane is for.
\end{attention}

\courssub{Saving a Presentation: Formats and Compatibility}

\textbf{Save} (\textbf{Ctrl+S}) updates the current file under its existing name; \textbf{Save As} (\textbf{F12}) creates a copy under a new name, location or format. Choose the format according to the purpose:

\begin{center}
\begin{tabularx}{\linewidth}{|l|X|X|}
\hline
\rowcolor{popblueL} \textbf{Format} & \textbf{Properties} & \textbf{When to use it} \\
\hline
\texttt{.pptx} & Editable PowerPoint format; keeps animations and fonts if available. & Working file; presenting from a computer with PowerPoint. \\
\hline
\texttt{.odp} & OpenDocument Presentation, native to LibreOffice Impress. & Working file in free-software environments. \\
\hline
\texttt{.pdf} & Fixed layout: looks identical everywhere, but not editable and loses animations. & Sending by email/WhatsApp, printing handouts, guaranteeing the display. \\
\hline
\texttt{.ppsx} & PowerPoint Show: opens directly in Slide Show mode, not in edit mode. & Distributing a ready-to-present file that launches immediately. \\
\hline
\end{tabularx}
\end{center}

\begin{attention}[Exam tip (GCE)]
Know the difference between \texttt{.pptx} (\textbf{editable}, requires compatible software and fonts) and \texttt{.pdf} (\textbf{fixed layout}, readable on any device but not modifiable, animations lost). Use \texttt{.pptx} while producing and presenting; export to \texttt{.pdf} to distribute or print. Also remember: a \texttt{.pptx} opened in Impress may shift slightly --- always test your file on the machine that will project it.
\end{attention}

\courssub{Planning Before Producing}

Professional presenters spend more time \emph{planning} than clicking. Before touching the software, answer: \emph{Who will listen? What must they remember? In what order?}

\begin{methode}[Planning a presentation in 5 steps]
\begin{enumerate}
\item \textbf{Know your audience}: parents, a jury, investors or classmates do not expect the same vocabulary, depth or tone.
\item \textbf{State one clear objective}: a single sentence, e.g.\ ``convince the jury that our water-filter project is feasible''.
\item \textbf{Build an outline}: introduction (hook, plan), body (2--4 main ideas), conclusion (summary, call to action).
\item \textbf{Limit the number of slides}: roughly one slide per minute of speaking; a 10-minute defence needs about 8--12 slides, not 40.
\item \textbf{Rehearse}: present aloud, time yourself, adjust.
\end{enumerate}
\end{methode}

\begin{popfigure}
\begin{tikzpicture}[node distance=0.55cm, every node/.style={font=\small},
box/.style={draw=popblue, fill=popblueL, rounded corners, minimum width=2.6cm, minimum height=0.9cm, align=center, text=popdark},
arr/.style={->, thick, poporange}]
\node[box, align=center] (a){\textbf{Audience}\\ who listens?};
\node[box, right=of a, align=center] (o){\textbf{Objective}\\ one clear goal};
\node[box, right=of o, align=center] (t){\textbf{Outline}\\ intro--body--end};
\node[box, right=of t, align=center] (s){\textbf{Slides}\\ few, visual};
\node[box, right=of s, align=center] (r){\textbf{Rehearsal}\\ timed practice};
\draw[arr] (a) -- (o);
\draw[arr] (o) -- (t);
\draw[arr] (t) -- (s);
\draw[arr] (s) -- (r);
\draw[arr, dashed] (r.south) to[bend right=25] node[below, font=\footnotesize, text=popmuted]{improve} (t.south);
\end{tikzpicture}
\legende{The presentation planning process: each step feeds the next, and rehearsal sends you back to improve the outline.}
\end{popfigure}

\begin{exoresolu}[Planning and producing a mini-presentation]
\textbf{Statement.} Your class must present ``The importance of mobile money in Cameroon'' in 5 minutes to new Lower Sixth students. (a) State the audience and one objective. (b) Propose an outline and a suitable number of slides. (c) Give the layout of slide 1 and of a body slide. (d) In which formats will you save, and why?
\tcblower
\textbf{Solution.}
\begin{itemize}
\item (a) \emph{Audience}: Lower Sixth students, beginners in ICT, so simple vocabulary and concrete examples (MTN MoMo, Orange Money). \emph{Objective}: ``Show that mobile money makes everyday payments faster and safer.''
\item (b) \emph{Outline}: Introduction (a market scene in Bamenda where a trader receives MoMo), Body (what mobile money is; advantages; precautions), Conclusion (summary + advice). At about one slide per minute: \textbf{6 slides} --- Title, Hook/photo, Definition, Advantages, Precautions, Conclusion.
\item (c) Slide 1: \emph{Title Slide} layout (title + names + class). A body slide such as ``Advantages'': \emph{Title and Content} layout with a bulleted list respecting the 6x6 rule, e.g.\ ``Fast transfers / No need to travel / Available 24/7 / Secure PIN''.
\item (d) Save the working file as \texttt{.pptx} (editable, to keep improving it and to project it), then export a \texttt{.pdf} copy to share on the class WhatsApp group so that everyone sees the same fixed layout even without PowerPoint.
\end{itemize}
\end{exoresolu}

\begin{exercice}[Your turn]
You must present your school to visiting students from another region in 8 minutes. (1) Define the audience and one objective. (2) Draw up an outline and state the ideal number of slides. (3) For each slide, indicate the layout you would choose. (4) Explain the difference between duplicating and hiding a slide, and give one situation where each is useful.
\end{exercice}

\begin{corrige}[Exercise]
(1) \emph{Audience}: visiting students of your age --- friendly tone, photos, little jargon. \emph{Objective}: ``Make visitors want to join our school.'' (2) Outline: Introduction (welcome, map), Body (history, results, clubs, facilities), Conclusion (invitation). About 8--10 slides. (3) Slide 1: Title Slide; welcome/map: Title and Content with picture; history: Title and Content (short bullets); clubs vs.\ results: Two Content to compare; final: Blank with a large group photo. (4) \emph{Duplicating} creates an identical copy of a slide --- useful to reuse a structure (e.g.\ one slide per club). \emph{Hiding} keeps a slide in the file but skips it during projection --- useful for a backup slide with extra statistics, shown only if a visitor asks a question.
\end{corrige}

\begin{aretenir}[Key points of Section 1]
\begin{itemize}
\item Presentation software (PowerPoint, Impress, Google Slides, WPS) produces a sequence of \cle{slides} that \emph{support} the speaker.
\item Know the window: ribbon, slide pane, thumbnail pane, notes pane, status bar, view buttons.
\item Choose \textbf{one} design theme; use \cle{layouts} and \cle{placeholders} instead of drawing text boxes everywhere.
\item Manage slides (add, duplicate, move, hide, delete) in the thumbnail pane or Slide Sorter view.
\item Respect the \cle{6x6 rule}; put details in the speaker notes.
\item Save as \texttt{.pptx} (editable) while working; export \texttt{.pdf} (fixed layout) to distribute; consider \texttt{.ppsx} for direct show mode.
\item Plan first: audience $\rightarrow$ objective $\rightarrow$ outline $\rightarrow$ slides $\rightarrow$ rehearsal.
\end{itemize}
\end{aretenir}

\coursec{Section 2: Designing Attractive and Consistent Slides}

In Section 1 we learned how to open presentation software, create slides, choose layouts and type content. A slide deck that merely \emph{contains} information, however, is not yet a good presentation. Think of two candidates presenting the same project on ``Mobile Money in Cameroon'' at a school science fair in Yaoundé: the first projects dense paragraphs in yellow text on a white background; the second shows a few large, well-aligned points in dark blue on white, with one clear picture per slide. Same content, completely different impact. This section teaches you the \cle{design principles} and the software tools --- themes, the \cle{Slide Master}, pictures, shapes, tables, charts, headers and footers --- that turn raw slides into an attractive, consistent and readable presentation.

\courssub{Design Principles: Consistency, Contrast, Alignment, Proximity, Simplicity}

Good slide design is not a matter of talent; it obeys a small number of rules that anyone can apply. Five principles, often summarised from graphic design theory, guide every decision you make on a slide.

\begin{definition}[The five design principles]
\begin{itemize}
\item \cle{Consistency}: the same fonts, colours, logo position and title style are used on \emph{every} slide, so the audience recognises the presentation as one coherent document.
\item \cle{Contrast}: important elements (titles, key figures) stand out through differences in colour, size or weight; text must contrast strongly with its background.
\item \cle{Alignment}: every object is lined up with something else (titles on the same vertical line, bullet blocks starting at the same left edge). Nothing is placed at random.
\item \cle{Proximity}: related items are grouped close together; unrelated items are separated by space, so the eye immediately sees what belongs together.
\item \cle{Simplicity}: one idea per slide, few words, generous empty space. A slide supports the speaker; it does not replace the speech.
\end{itemize}
\end{definition}

\begin{exemplebox}[Choosing fonts and colours]
A student preparing a talk for the GCE ICT practical chooses:
\begin{itemize}
\item \textbf{Fonts}: one sans-serif font for titles (e.g.\ Calibri or Arial, bold) and the same family for body text --- never more than \textbf{two or three fonts} in a whole presentation.
\item \textbf{Font sizes}: titles at least $36$--$44$ points, body text at least $24$ points, so that text remains readable from the back of a classroom.
\item \textbf{Colours}: a dark text colour (navy blue or black) on a light background (white or cream), plus one accent colour (e.g.\ orange) used only for emphasis.
\end{itemize}
\end{exemplebox}

\begin{attention}[Common mistakes in design]
The most frequent errors are: using \textbf{light text on a light background} (yellow on white, light grey on beige); using \textbf{more than three different fonts}; and filling every centimetre of the slide with text and clip art. Remember that classroom projectors \emph{wash out} low-contrast colours: what looks acceptable on your laptop screen can become invisible on the wall of a bright classroom in Buea or Bamenda. Always test with the projector if possible.
\end{attention}

\begin{popfigure}
\begin{tikzpicture}
% Cluttered slide
\draw[fill=popblueL, draw=popmuted] (0,0) rectangle (5.6,3.8);
\node[anchor=west, font=\scriptsize\bfseries, text=poppurple] at (0.2,3.5) {Mobile money!!!};
\node[anchor=west, font=\tiny, text=popdark, align=left, text width=5cm] at (0.25,2.6) {Mobile money is very important\\ in Cameroon because many\\ people use MTN MoMo and\\ Orange Money every day to\\ send and receive money and\\ pay bills and buy airtime...};
\node[font=\tiny, text=popgold] at (4.3,3.3) {CLICK HERE};
\draw[fill=poppink, draw=popdark] (3.6,0.4) rectangle (5.3,1.6);
\node[font=\tiny, text=white] at (4.45,1.0) {photo};
\draw[fill=popgreenL, draw=popdark] (0.3,0.3) circle (0.35);
\node[font=\tiny] at (0.65,0.3) {art};
\node[font=\scriptsize, text=popmuted] at (2.8,-0.4) {Cluttered: many fonts, low contrast, no alignment};
% Clean slide
\draw[fill=white, draw=popblue, thick] (6.4,0) rectangle (12,3.8);
\node[anchor=west, font=\small\bfseries, text=popblue] at (6.7,3.3) {Mobile Money in Cameroon};
\draw[poporange, thick] (6.7,3.05) -- (9.2,3.05);
\node[anchor=west, font=\scriptsize, text=popdark, align=left, text width=5cm] at (6.7,2.4) {\textbullet\ Send and receive money\\[2pt] \textbullet\ Pay bills and airtime\\[2pt] \textbullet\ Grow small businesses};
\draw[fill=popblueL, draw=popblue] (9.9,1.0) rectangle (11.7,2.6);
\node[font=\tiny, text=popblue] at (10.8,1.8) {one image};
\node[font=\scriptsize, text=popmuted] at (9.2,-0.4) {Clean: strong contrast, aligned, one idea};
\end{tikzpicture}
\legende{A cluttered slide (left) versus a clean slide (right) respecting contrast, alignment and simplicity.}
\end{popfigure}

\courssub{Themes, Variants and Background Styles}

Applying the five principles slide by slide would be slow and error-prone. Presentation software solves this with \cle{themes}.

\begin{definition}[Theme and variant]
A \cle{theme} is a coordinated set of colours, fonts and effects that is applied to the whole presentation at once (tab \textbf{Design} $>$ \textbf{Themes}). A \cle{variant} is an alternative colour or font combination \emph{within} the same theme, offered in the \textbf{Variants} group of the Design tab.
\end{definition}

\begin{methode}[Applying a theme and changing the background]
\begin{enumerate}
\item Open the \textbf{Design} tab and click a theme thumbnail: all slides instantly adopt its colours and fonts.
\item In the \textbf{Variants} group, click a colour variant to switch, for example, from a blue to a green palette while keeping the same fonts.
\item To personalise further, click \textbf{Format Background}: choose a solid fill, gradient, picture or texture, then click \textbf{Apply to All} so that every slide shares the same background.
\end{enumerate}
\end{methode}

\begin{attention}[Backgrounds and readability]
A busy background picture behind text destroys contrast. If you use a picture background, darken it or place a semi-transparent solid shape behind the text. When in doubt, a plain light background is always the safest choice for projection.
\end{attention}

\courssub{The Slide Master: One Change, Every Slide}

Suppose your school asks you to place its logo and the slide number on \emph{all 30 slides} of your presentation. Copying and pasting the logo 30 times is slow, and if the logo changes you must redo everything. The correct tool is the \cle{Slide Master}.

\begin{definition}[Slide Master]
The \cle{Slide Master} is the top-level slide that stores the theme, fonts, colours and objects (logo, footer, date, slide number) inherited by all the slides that use its \cle{layouts}. It is opened from \textbf{View} $>$ \textbf{Slide Master}.
\end{definition}

\begin{propriete}[Inheritance from the Slide Master]
Any change made on the Slide Master --- inserting a logo, setting the title font, adding a footer --- automatically appears on \textbf{every slide} using that master. Similarly, a change made on a \emph{layout} (e.g.\ the ``Title and Content'' layout) appears on all slides using that layout only. (The mechanism is examinable at the GCE; memorise the hierarchy Master $\to$ Layouts $\to$ Slides.)
\end{propriete}

\begin{popfigure}
\begin{tikzpicture}[
 master/.style={rectangle, draw=popblue, fill=popblueL, thick, minimum width=3.2cm, minimum height=1cm, font=\small\bfseries, text=popdark},
 lay/.style={rectangle, draw=poporange, fill=white, thick, minimum width=2.6cm, minimum height=0.8cm, font=\scriptsize, text=popdark},
 sli/.style={rectangle, draw=popgreen, fill=popgreenL, minimum width=2.2cm, minimum height=0.7cm, font=\scriptsize, text=popdark},
 arr/.style={->, thick, popmuted}]
\node[master] (m) at (4.5,5) {Slide Master};
\node[font=\tiny, text=popmuted] at (4.5,4.35) {theme, fonts, logo, footer};
\node[lay] (l1) at (1.6,3) {Title layout};
\node[lay] (l2) at (4.5,3) {Title \& Content};
\node[lay] (l3) at (7.4,3) {Two Content};
\node[sli] (s1) at (1.6,1.2) {Slide 1};
\node[sli] (s2) at (3.7,1.2) {Slide 2};
\node[sli] (s3) at (5.5,1.2) {Slide 3};
\node[sli] (s4) at (7.4,1.2) {Slide 4};
\draw[arr] (m) -- (l1);
\draw[arr] (m) -- (l2);
\draw[arr] (m) -- (l3);
\draw[arr] (l1) -- (s1);
\draw[arr] (l2) -- (s2);
\draw[arr] (l2) -- (s3);
\draw[arr] (l3) -- (s4);
\end{tikzpicture}
\legende{Hierarchy of inheritance: formatting placed on the Slide Master flows down through the layouts to every slide.}
\end{popfigure}

\begin{methode}[Adding a logo and slide number to all slides]
\begin{enumerate}
\item Click \textbf{View} $>$ \textbf{Slide Master}.
\item Select the top (largest) master slide in the left pane.
\item Insert the logo (\textbf{Insert} $>$ \textbf{Pictures}) and position it, e.g.\ bottom-right corner.
\item Use \textbf{Insert} $>$ \textbf{Header \& Footer} to tick \emph{Slide number} and \emph{Date and time}, then \textbf{Apply to All}.
\item Click \textbf{Close Master View}: the logo and numbers now appear on every slide.
\end{enumerate}
\end{methode}

\begin{exemplebox}[GCE exam tip]
Exam questions very often ask: ``How can the same logo (or slide number) be added to all the slides of a presentation at once?'' The expected answer is: \textbf{use the Slide Master} (View $>$ Slide Master), \emph{not} copying the object onto each slide one by one. Copying is accepted only as a wrong-method contrast in your explanation.
\end{exemplebox}

\courssub{Inserting and Formatting Pictures}

Pictures make a slide attractive and memorable --- a photo of a mobile money kiosk says more than a paragraph. Pictures can be inserted from a \textbf{file} on the computer (\textbf{Insert} $>$ \textbf{Pictures} $>$ \textbf{This Device}) or from \textbf{online sources} (Insert $>$ Online Pictures / Stock Images).

Once inserted, a picture can be formatted with the \textbf{Picture Format} tab:
\begin{itemize}
\item \cle{Cropping}: cutting away unwanted parts of the image to focus on the subject.
\item \cle{Resizing}: dragging a \emph{corner} handle keeps the proportions; dragging a side handle distorts the image.
\item \cle{Borders and effects}: frames, shadows, reflections and soft edges from the Picture Styles gallery.
\item \cle{Compressing images}: \textbf{Picture Format} $>$ \textbf{Compress Pictures} reduces file size (essential when the presentation must fit on a flash drive or be sent by email), at the cost of some quality.
\end{itemize}

\begin{attention}[Stretching and copyright]
Never stretch a picture sideways: faces and objects become deformed. Always resize from a corner handle. Also, images taken from the Internet may be protected by copyright; for a school presentation, prefer the built-in stock images or free-image sites, and mention the source.
\end{attention}

\courssub{Shapes, Text Boxes, SmartArt and Icons}

Beyond pictures, slides use drawn objects to explain ideas:
\begin{itemize}
\item \cle{Shapes} (Insert $>$ Shapes): arrows, rectangles, callouts --- ideal for simple diagrams such as a flow of money between customer, agent and bank.
\item \cle{Text boxes} (Insert $>$ Text Box): place text anywhere, outside the predefined placeholders.
\item \cle{SmartArt} (Insert $>$ SmartArt): ready-made diagrams --- lists, processes, cycles, hierarchies --- that convert bullet points into an attractive graphic in one click.
\item \cle{Icons} (Insert $>$ Icons): small standard symbols (phone, money, shield) that illustrate a point without a photo.
\end{itemize}

When several objects form one figure, select them all (hold \textbf{Shift} and click, or drag a selection box), then:
\begin{itemize}
\item \textbf{Group} (right-click $>$ Group) so they move and resize as a single object;
\item \textbf{Align} (Shape Format $>$ Align) to line them up (Align Left, Align Top, Distribute Horizontally) --- this is how the \emph{alignment} principle is applied in practice.
\end{itemize}

\begin{experience}[Guided practice: Building a process diagram with SmartArt]
\textbf{Goal:} Create a slide that explains the steps of an MTN Mobile Money cash withdrawal.
\begin{enumerate}
\item Insert a new slide with \textbf{Title and Content} layout. Title: ``Withdrawing Cash via MTN MoMo''.
\item Click the \textbf{SmartArt} icon in the content placeholder, choose \textbf{Process} $>$ \textbf{Basic Chevron Process}, click OK.
\item In the text pane, type the four steps: \textbullet\ Dial *150\# \textbullet\ Select ``Withdraw Cash'' \textbullet\ Enter amount and agent code \textbullet\ Receive cash from agent.
\item With the SmartArt selected, use \textbf{SmartArt Design} $>$ \textbf{Change Colors} to pick a colour scheme matching your theme.
\item Use \textbf{Format} $>$ \textbf{Size} to make the diagram large enough to be read from the back of the room.
\end{enumerate}
\textbf{Observation:} The diagram stays perfectly aligned; if you add a step, the layout adjusts automatically.
\end{experience}

\courssub{Tables and Charts on a Slide}

Numerical information is clearer as a table or a chart than as sentences.

\begin{definition}[Table and chart on a slide]
A \cle{table} organises data in rows and columns (Insert $>$ Table); it is formatted with the Table Design tab (header shading, borders, banded rows). A \cle{chart} (Insert $>$ Chart) represents data graphically: column charts compare values, pie charts show proportions of a whole, line charts show evolution over time.
\end{definition}

\begin{methode}[Inserting a chart]
\begin{enumerate}
\item Click \textbf{Insert} $>$ \textbf{Chart}.
\item Choose the type: \textbf{column} (comparison), \textbf{pie} (shares) or \textbf{line} (trend), then OK.
\item A small spreadsheet (data sheet) opens: replace the sample data with your own figures; the chart updates automatically.
\item Close the data sheet, then use the \textbf{Chart Design} and \textbf{Format} tabs to add a title, data labels and a legend.
\end{enumerate}
\end{methode}

\begin{propriete}[Linking chart data to a spreadsheet]
The chart on a slide is linked to its data sheet: editing a number in the sheet immediately updates the chart. A chart can also be \emph{linked} to an external spreadsheet file, so that when the spreadsheet is updated, the presentation chart can be refreshed too (Paste Special $>$ Paste Link when copying a chart from a spreadsheet).
\end{propriete}

\begin{exoresolu}[Choosing and building the right chart]
\textbf{Statement.} A Lower Sixth student surveyed how 40 classmates pay their canteen expenses: 22 use mobile money, 12 use cash, 6 use a bank card. State the most suitable chart type and describe how to insert it.
\tcblower
\textbf{Solution.} The data are \emph{shares of a whole} (the 40 students), so a \textbf{pie chart} is the most suitable type (a column chart would also be acceptable for comparison, but a line chart would be wrong: there is no time evolution).
Steps: (1) Insert $>$ Chart $>$ Pie $>$ OK. (2) In the data sheet that opens, enter the three categories --- Mobile money 22, Cash 12, Bank card 6. (3) Close the sheet; the pie slices appear automatically in proportion. (4) Add data labels showing percentages ($55\%$, $30\%$, $15\%$) and a title ``Payment methods in the canteen''. (5) Format with strong contrasting slice colours so the chart stays readable on a projector.
\end{exoresolu}

\courssub{Headers, Footers, Slide Numbers and Speaker Notes}

Professional slides carry small pieces of repeated information:
\begin{itemize}
\item \cle{Header / Footer}: a line of text (e.g.\ the school name or presentation title) repeated on slides, set via \textbf{Insert} $>$ \textbf{Header \& Footer}.
\item \cle{Slide number}: automatic numbering, also from Header \& Footer; it helps the audience refer to a slide during questions.
\item \cle{Date}: fixed or updated automatically each time the file is opened.
\end{itemize}
As seen above, placing these on the Slide Master guarantees they appear everywhere consistently.

\begin{definition}[Speaker notes]
\cle{Speaker notes} are private notes typed in the pane below the slide (or in View $>$ Notes Page). They are visible to the presenter (in Presenter View or on printed notes pages) but \textbf{not} projected to the audience. They let the speaker keep the slide simple while remembering details, figures and examples.
\end{definition}

\begin{exemplebox}[Speaker notes in practice]
On a slide showing only the pie chart of canteen payments, the presenter types in the notes pane: ``Explain that mobile money dominates because an MTN kiosk stands just in front of the school gate; mention that card payment requires a POS terminal the canteen does not have.'' The slide stays clean; the speech stays rich.
\end{exemplebox}

\courssub{Checking Readability Before Presenting}

Before delivering the presentation, always perform a final readability check, ideally with the actual projector:

\begin{aretenir}[Readability checklist]
\begin{itemize}
\item \textbf{Font sizes}: titles $\geq 36$ pt, body text $\geq 24$ pt; if you must shrink text to fit, you have too much text.
\item \textbf{Contrast}: dark text on light background (or the reverse); avoid red/green combinations and pale colours that projectors wash out in bright classrooms.
\item \textbf{Clutter}: maximum about 6 lines of text per slide and 6 words per line (the ``6 $\times$ 6'' rule of thumb); one idea per slide.
\item \textbf{Consistency}: same theme, same fonts, logo and slide numbers on every slide --- check via the Slide Master.
\item \textbf{Proofreading}: run the spelling checker and read each slide aloud once.
\end{itemize}
\end{aretenir}

\begin{exercice}[Designing a consistent presentation]
Your school asks you to prepare a 10-slide presentation on ``The use of ICT in Cameroonian hospitals''.
\begin{enumerate}
\item Explain how you would guarantee that the hospital logo, the slide number and the same fonts appear on all 10 slides without editing them one by one.
\item You have a table of the number of connected hospitals per region for the years 2021, 2022 and 2023. Which chart type best shows the \emph{evolution} over time, and how do you insert it?
\item State two readability checks you would perform before projecting in a bright hall.
\end{enumerate}
\end{exercice}

\begin{corrige}[Exercise 1]
\begin{enumerate}
\item Open \textbf{View} $>$ \textbf{Slide Master}, select the top master, insert the logo and position it, set the title and body fonts, then use \textbf{Insert} $>$ \textbf{Header \& Footer} to tick \emph{Slide number} and \textbf{Apply to All}. Every slide inherits these elements automatically.
\item A \textbf{line chart} (one line per region, years on the horizontal axis) best shows evolution over time. Insert it via \textbf{Insert} $>$ \textbf{Chart} $>$ \textbf{Line}, then type the years and values in the data sheet that opens; the chart updates automatically.
\item Check that body text is at least 24 pt and titles at least 36 pt, and that there is strong contrast (dark text on a light background), because projectors in bright rooms wash out pale colours; also remove any clutter so each slide carries one idea only.
\end{enumerate}
\end{corrige}

\begin{savaistu}[Did you know?]
In Cameroon, public services and companies such as CAMTEL, MTN and Orange use standardised presentation templates built on slide masters: the same logo, colours and footer appear on every slide of every official presentation. This is exactly the \emph{consistency} principle --- and it is why mastering the Slide Master is a real professional skill, not just an exam topic.
\end{savaistu}

With well-designed, consistent and readable slides in hand, the next step is to bring them to life: Section 3 will add multimedia --- sound, video --- together with transitions between slides and animations of objects on a slide.

\coursec{Section 3: Multimedia, Transitions and Animations}

In the previous section we learned how to design visually consistent slides using themes, slide masters, and layout principles. A well‑designed static slide is a solid foundation, but modern presentations often need to \emph{show} processes, \emph{play} media, and \emph{guide} the audience’s attention dynamically. This section introduces the tools that bring slides to life: multimedia insertion, slide transitions, object animations, and interactive navigation. Mastering these features will not only make your presentations more engaging but also earn you full marks in GCE practical tasks where examiners explicitly test the distinction between transitions and animations, the correct use of the Animation Pane, and the ability to create a self‑running kiosk‑style show.

\courssub{Multimedia in Presentations}

Modern presentation software (Microsoft PowerPoint, LibreOffice Impress, Google Slides) treats audio and video as first‑class objects that can be placed on a slide just like an image. The most widely supported formats in the Cameroonian school and office environment are:

\begin{itemize}
\item \textbf{Audio:} \cle{MP3} (compressed, small size), \cle{WAV} (uncompressed, high quality), and increasingly \cle{M4A}/AAC.
\item \textbf{Video:} \cle{MP4} with H.264 video codec and AAC audio — the de facto standard for compatibility across Windows, macOS, Linux, and mobile devices.
\end{itemize}

\begin{definition}[Multimedia Object]
A multimedia object is an audio or video file embedded in or linked to a slide. When embedded, the file becomes part of the presentation package (larger file size, guaranteed playback). When linked, only the path is stored (smaller file, but the media file must travel with the presentation).
\end{definition}

\textbf{Inserting media.}\par
Use \texttt{Insert \$\textbackslash{}to\$ Media \$\textbackslash{}to\$ Audio/Video \$\textbackslash{}to\$ This Device...} (PowerPoint) or \texttt{Insert \$\textbackslash{}to\$ Audio or Video} (Impress). The media appears as a speaker icon (audio) or a video frame (video) that you can resize and reposition.

\textbf{Playback options.}\par
Select the media object to reveal the \texttt{Playback} (or \texttt{Media Options}) tab. Key settings:
\begin{itemize}
\item \textbf{Start:} \cle{Automatically} (plays as soon as the slide appears), \cle{On Click} (waits for a mouse/key press), or \cle{In Click Sequence} (part of the animation order).
\item \textbf{Loop until Stopped:} repeats the media continuously — useful for background music on a title slide.
\item \textbf{Trim Media:} lets you set start and end points without an external editor. Click \texttt{Trim}, drag the green/red markers, and confirm.
\item \textbf{Fade In/Out:} adds smooth volume ramps.
\item \textbf{Hide During Show:} (audio only) hides the speaker icon during the slide show.
\item \textbf{Play Across Slides:} (audio only) continues playback across subsequent slides.
\end{itemize}

\begin{exemplebox}[Adding a video clip to a slide about Mount Cameroon]
You are preparing a slide on tourism in the Southwest Region. You have a 2‑minute MP4 clip of the Mount Cameroon Race of Hope. Insert the video, resize it to occupy the right half of the slide, set \textbf{Start: On Click} so you can introduce it verbally, and use \textbf{Trim} to keep only the 15‑second summit sprint. Tick \textbf{Hide While Not Playing} if you want the frame invisible before playback.
\end{exemplebox}

\begin{attention}[Common Mistake — Missing Media Files]
If you choose \emph{Link to File} instead of \emph{Insert} (embed), the video will not play on another computer unless you copy the media file into the same folder as the presentation. For GCE practical exams, always \textbf{embed} media unless the question explicitly asks for linking.
\end{attention}

\courssub{Recording Narration and Timings for a Self‑Running Presentation}

A \cle{self‑running presentation} (kiosk mode) advances automatically and includes a recorded voice‑over. This is ideal for unattended displays at school open days, CAMTEL customer‑care screens, or GCE practical submissions where the examiner watches the show without your live narration.

\begin{methode}[Recording a narrated, self‑running show]
\begin{enumerate}
\item Open the presentation and go to \texttt{Slide Show \$\textbackslash{}to\$ Record Slide Show} (PowerPoint) or \texttt{Slide Show \$\textbackslash{}to\$ Record} (Impress).
\item Choose \texttt{Record from Beginning}. A recording toolbar appears with \texttt{Record}, \texttt{Stop}, \texttt{Replay}, and navigation arrows.
\item Speak clearly into your microphone. Click \texttt{Next} (or press \texttt{Space}) to advance slides; the time spent on each slide is recorded as \textbf{slide timings}.
\item Use the \texttt{Pen} or \texttt{Highlighter} tools during recording to annotate slides — these strokes become part of the show.
\item Press \texttt{Stop} when finished. Save the file; the timings and narration are now embedded.
\item To enforce kiosk mode: \texttt{Slide Show \$\textbackslash{}to\$ Set Up Slide Show \$\textbackslash{}to\$ Browsed at a kiosk (full screen)}. The show will now loop until \texttt{Esc} is pressed.
\end{enumerate}
\end{methode}

\begin{experience}[Guided Practice — Create a 3‑Slide Auto‑Advancing Show]
\begin{enumerate}
\item Create three slides: Title, “MTN Mobile Money Steps”, “Orange Money Steps”.
\item Insert a 5‑second MP3 jingle on the title slide; set \textbf{Start: Automatically}, \textbf{Loop until Stopped}.
\item Record a 10‑second narration on each content slide describing the USSD codes (*126\# for MTN, *150\# for Orange).
\item Set \texttt{Slide Show \$\textbackslash{}to\$ Set Up Slide Show \$\textbackslash{}to\$ Browsed at a kiosk}.
\item Press \texttt{F5} and verify the show runs unattended for 30 seconds, then loops.
\end{enumerate}
\end{experience}

\courssub{Slide Transitions}

\begin{definition}[Transition]
A \cle{transition} is the visual effect that plays when moving from one slide to the next. It belongs to the \emph{slide} as a whole, not to individual objects.
\end{definition}

Transitions are managed on the \texttt{Transitions} tab. Common built‑in effects include:
\begin{itemize}
\item \textbf{Fade:} smooth cross‑dissolve — professional and subtle.
\item \textbf{Push:} new slide pushes the old one out — conveys progression.
\item \textbf{Wipe:} a line sweeps across the screen — directional.
\item \textbf{Morph:} (PowerPoint 2019+/365) automatically animates differences between two similar slides — powerful for object movement, resizing, and colour changes without manual animation.
\item \textbf{Cut:} instantaneous change — no visual effect.
\end{itemize}

\textbf{Configuring a transition.}\par
Select one or more slides (Ctrl+Click for multiple, or Ctrl+A for all). Choose an effect from the gallery. Then adjust:
\begin{itemize}
\item \textbf{Effect Options:} direction (e.g., Push from Left, Right, Top, Bottom).
\item \textbf{Duration:} how long the transition lasts (default 1.00 s; 0.50–1.50 s is typical).
\item \textbf{Sound:} optional chime, whoosh, or custom WAV — use sparingly.
\item \textbf{Advance Slide:}
      \begin{itemize}
      \item \cle{On Mouse Click} (default) — presenter controls pace.
      \item \cle{After:} \SI{5}{s} (or any time) — slide advances automatically after the set delay. Combined with recorded timings, this creates a self‑running show.
      \end{itemize}
\end{itemize}

\begin{aretenir}[Transition Best Practices]
\begin{itemize}
\item Apply \textbf{one consistent transition} (e.g., Fade) to the whole presentation via \texttt{Apply to All}.
\item Reserve distinct transitions for structural breaks (e.g., a \textbf{Push} when moving from “Problem” to “Solution”).
\item Keep \textbf{Duration} between 0.5 s and 1.5 s; longer transitions test the audience’s patience.
\item Avoid sounds unless the presentation runs unattended in a noisy environment (e.g., a trade fair booth at the Yaoundé Congress Centre).
\item For GCE practicals, if the question says “apply a transition to all slides”, use \texttt{Apply to All} after configuring the effect, duration, and advance settings.
\end{itemize}
\end{aretenir}

\begin{attention}[GCE Trap — Transition vs. Animation]
Examiners frequently ask: “State the difference between a slide transition and an animation.” The correct answer: \textbf{A transition occurs between slides; an animation occurs on objects within a slide.} Confusing the two loses marks instantly.
\end{attention}

\courssub{Object Animations}

\begin{definition}[Animation]
An \cle{animation} is a movement or effect applied to an individual object (text box, picture, shape, chart, SmartArt) \emph{within a single slide}. It controls how the object appears, behaves, or disappears.
\end{definition}

Animations are organised into four categories. The diagram below summarises them with a representative icon for each.

\begin{popfigure}
\begin{tikzpicture}[
  cat/.style={draw=popblue, thick, rounded corners, fill=popblueL, minimum width=3.2cm, minimum height=2.8cm, text width=3cm, align=center, font=\small},
  icon/.style={font=\Large, text=popdark},
  arr/.style={->, thick, popline}
]
\node[cat, align=center] (ent) at (0,0){\textbf{Entrance}\\\vspace{0.2cm}\tikz[baseline=-0.5ex]\node[icon]{$\blacktriangleright$}; Fade In, Fly In, Zoom};
\node[cat, align=center] (emp) at (4.5,0){\textbf{Emphasis}\\\vspace{0.2cm}\tikz[baseline=-0.5ex]\node[icon]{$\ast$}; Pulse, Grow/Shrink, Spin};
\node[cat, align=center] (ext) at (9,0){\textbf{Exit}\\\vspace{0.2cm}\tikz[baseline=-0.5ex]\node[icon]{$\blacktriangleleft$}; Fade Out, Fly Out, Disappear};
\node[cat, align=center] (mot) at (13.5,0){\textbf{Motion Path}\\\vspace{0.2cm}\tikz[baseline=-0.5ex]\node[icon]{$\rightsquigarrow$}; Line, Arc, Custom Path};

\draw[arr] (ent.north east) -- (emp.north west) node[midway, above, font=\tiny] {same object can have multiple};
\draw[arr] (emp.north east) -- (ext.north west) node[midway, above, font=\tiny] {sequence: enter $\to$ emphasize $\to$ exit};
\draw[arr] (ext.north east) -- (mot.north west) node[midway, above, font=\tiny] {motion path can be any category};
\end{tikzpicture}
\legende{The four animation categories with a typical example icon. An object can receive multiple animations (e.g., Entrance $\to$ Emphasis $\to$ Exit).}
\end{popfigure}

\textbf{Applying an animation.}\par
Select the object $\to$ \texttt{Animations} tab $\to$ open the gallery (click the down arrow) $\to$ choose an effect. A small numbered tag appears next to the object indicating its order in the animation sequence.

\textbf{Preview.}\par
Click \texttt{Preview} on the Animations tab or press \texttt{Shift+F5} to start the show from the current slide.

\begin{exemplebox}[Animating a photo of the Reunification Monument]
Insert a picture of the monument in Buea. Apply \textbf{Entrance: Fade} (duration 1 s). Then \textbf{Add Animation} $\to$ \textbf{Emphasis: Grow/Shrink} (1.5 s) to draw attention. Finally \textbf{Add Animation} $\to$ \textbf{Exit: Fade} (0.5 s) to clear the slide for the next point. The tags will show 1, 2, 3.
\end{exemplebox}

\courssub{The Animation Pane and Sequencing}

The \cle{Animation Pane} (PowerPoint) or \cle{Custom Animation} pane (Impress) is the control centre for fine‑tuning animation behaviour. Open it via \texttt{Animations \$\textbackslash{}to\$ Animation Pane}.

For each effect you can set:
\begin{itemize}
\item \textbf{Start:}
      \begin{itemize}
      \item \cle{On Click} — waits for the presenter’s click (default).
      \item \cle{With Previous} — starts simultaneously with the previous effect in the list.
      \item \cle{After Previous} — starts automatically after the previous effect finishes.
      \end{itemize}
\item \textbf{Duration:} length of the effect in seconds (two decimals).
\item \textbf{Delay:} pause before the effect starts (useful for staggering).
\item \textbf{Reorder:} drag effects up/down or use the \texttt{Reorder Animation} arrows.
\item \textbf{Effect Options:} direction, sequence (by paragraph, by series, by category, etc.), and other effect‑specific settings.
\end{itemize}

The timeline below illustrates how \texttt{Start} options affect the sequence of three effects (A, B, C) on the same slide.

\begin{popfigure}
\begin{tikzpicture}[
  bar/.style={draw=popdark, thick, fill=popblueL, rounded corners=2pt},
  lab/.style={font=\small, text=popdark},
  time/.style={font=\tiny, text=popmuted}
]
% Time axis
\draw[thick, ->] (0,0) -- (14,0) node[right, lab] {Time (seconds)};
\foreach \x/\t in {0/0, 2/2, 4/4, 6/6, 8/8, 10/10, 12/12}
  \draw (\x, -0.1) -- (\x, 0.1) node[below, time] {\t};

% Effect A: On Click at t=0, duration 2s
\draw[bar] (0, 2.5) rectangle (4, 3.0) node[midway, lab] {A: On Click (2s)};
\draw[dashed, popmuted] (0,2.5) -- (0,0);

% Effect B: With Previous (starts with A), duration 3s
\draw[bar, fill=popgreenL] (0, 1.8) rectangle (6, 2.3) node[midway, lab] {B: With Previous (3s)};

% Effect C: After Previous (after B ends at 6s), duration 1.5s
\draw[bar, fill=poporange] (6, 1.1) rectangle (9, 1.6) node[midway, lab] {C: After Previous (1.5s)};
\draw[dashed, popmuted] (6,1.1) -- (6,0);

% Labels for start types
\node[left, lab] at (-0.5, 2.75) {\textbf{On Click}};
\node[left, lab] at (-0.5, 2.05) {\textbf{With Previous}};
\node[left, lab] at (-0.5, 1.35) {\textbf{After Previous}};
\end{tikzpicture}
\legende{Animation timeline showing three effects with different start options. Effect A starts on click at t=0. Effect B starts with A (simultaneous). Effect C starts after B finishes (t=6 s).}
\end{popfigure}

\begin{methode}[Animating a bulleted list progressively (progressive disclosure)]
\begin{enumerate}
\item Select the text box containing the bulleted list.
\item On the \texttt{Animations} tab, choose an \textbf{Entrance} effect (e.g., \textbf{Fade} or \textbf{Appear}).
\item Open the \texttt{Animation Pane}. The list appears as a single effect with a double‑arrow icon indicating “grouped by paragraph”.
\item Click the effect’s drop‑down arrow $\to$ \textbf{Effect Options} $\to$ \textbf{Text Animation} tab.
\item Set \textbf{Group text:} \cle{By 1st Level Paragraphs} (or \cle{By 2nd Level} for sub‑bullets).
\item In the Animation Pane, each bullet now appears as a separate sub‑effect. Set \textbf{Start: On Click} for each (default).
\item Preview: each click reveals the next bullet, keeping the audience focused on the current point.
\end{enumerate}
\end{methode}

\begin{exemplebox}[Animating a column chart by series]
Insert a clustered column chart comparing MTN vs. Orange market share (2020–2023). Apply \textbf{Entrance: Wipe} (From Bottom). In \texttt{Effect Options \$\textbackslash{}to\$ Chart Animation}, choose \textbf{By Series}. Each year’s pair of columns appears together on click, allowing you to comment year by year.
\end{exemplebox}

\begin{exemplebox}[Animating a SmartArt diagram]
Insert a \texttt{Process} SmartArt showing the steps of Mobile Money registration. Apply \textbf{Entrance: Fade} and set \textbf{Effect Options $\to$ One by One}. Each shape appears on click, letting you explain each step before revealing the next.
\end{exemplebox}

\courssub{Hyperlinks and Action Buttons}

Interactive presentations let the viewer (or presenter) jump to specific slides, external files, or web pages. This is essential for non‑linear shows like a digital kiosk at the National Museum or a revision menu for GCE candidates.

\textbf{Hyperlinks.}\par
Select text, a shape, or an image $\to$ \texttt{Insert \$\textbackslash{}to\$ Link} (or \texttt{Ctrl+K}). In the dialog:
\begin{itemize}
\item \textbf{Place in This Document:} choose a slide title (e.g., “Question 3”).
\item \textbf{Existing File or Web Page:} paste a URL (e.g., \texttt{https://www.antic.cm} for ANTIC’s cybersecurity portal) or browse to a local PDF.
\item \textbf{Email Address:} creates a \texttt{mailto:} link.
\item \textbf{ScreenTip:} add a tooltip that appears on hover — good for accessibility.
\end{itemize}

\textbf{Action Buttons.}\par
\texttt{Insert \$\textbackslash{}to\$ Shapes \$\textbackslash{}to\$ Action Buttons} (bottom of the gallery). Draw the button on the slide; a dialog appears instantly to assign:
\begin{itemize}
\item \textbf{Hyperlink to:} Next Slide, Previous Slide, First Slide, Last Slide, Last Slide Viewed, Specific Slide, URL, Other File.
\item \textbf{Run Program:} launch an external .exe (rarely used in exams).
\item \textbf{Play Sound:} optional click feedback.
\item \textbf{Mouse Over:} (in the Action Settings dialog) trigger on hover instead of click.
\end{itemize}

\begin{experience}[Build an Interactive Revision Menu]
Create a slide titled “GCE ICT Revision”. Add four Action Buttons labelled “Hardware”, “Software”, “Networks”, “Databases”. Hyperlink each to the corresponding section slide. On each section slide, place a small “Home” action button (House icon) linking back to the menu. Test in Slide Show mode: clicking jumps instantly without linear navigation.
\end{experience}

\courssub{Good Practice and Exam Tips}

\begin{aretenir}[Design Principles for Transitions and Animations]
\begin{enumerate}
\item \textbf{Substance over spectacle:} Every effect should clarify the message. A \textbf{Fade} transition signals a new topic; a \textbf{Motion Path} shows a process flow (e.g., data moving from CPU to RAM).
\item \textbf{Consistency:} Use one transition style and one entrance animation family throughout.
\item \textbf{Speed:} Fast enough to feel responsive (0.5–1 s), slow enough to be perceived.
\item \textbf{Sound:} Only for self‑running kiosks or accessibility cues; never for live presentations.
\item \textbf{Accessibility:} Avoid flashing effects (>3 Hz) that can trigger photosensitive epilepsy. Provide alternative text for animated charts.
\item \textbf{Test on target hardware:} A presentation that runs smoothly on a modern laptop may lag on an older school computer. Keep file sizes reasonable; compress media if needed (\texttt{File \$\textbackslash{}to\$ Info \$\textbackslash{}to\$ Media Size and Performance}).
\end{enumerate}
\end{aretenir}

\begin{attention}[Common Mistake — Overloading with Effects]
Students often apply \emph{Spin}, \emph{Bounce}, and \emph{Sound} to every bullet. Examiners penalise this as “distracting use of animations”. In the GCE practical, if the marking scheme says “apply appropriate animations”, a simple \textbf{Fade} or \textbf{Appear} on click scores full marks; a spinning, sound‑emitting bullet may lose the “appropriateness” mark.
\end{attention}

\begin{attention}[Common Mistake — Forgetting to Embed Fonts]
If you use non‑standard fonts (e.g., downloaded from the internet), the presentation may look different on the examiner’s computer. Always \texttt{File \$\textbackslash{}to\$ Options \$\textbackslash{}to\$ Save \$\textbackslash{}to\$ Embed fonts in the file} (PowerPoint) or ensure standard fonts (Arial, Calibri, Times New Roman) are used.
\end{attention}

\begin{propriete}[GCE Exam Tip — Self‑Running Show Setup]
To make a slideshow advance automatically without any clicks:
\begin{enumerate}
\item Go to \texttt{Transitions} tab.
\item Uncheck \textbf{On Mouse Click}, check \textbf{After:} and set a time (e.g., \SI{10}{s}).
\item Click \texttt{Apply to All} (or set per slide).
\item Optionally, use \texttt{Slide Show \$\textbackslash{}to\$ Rehearse Timings} to record exact durations per slide while practising your narration.
\item Finally, \texttt{Set Up Slide Show \$\textbackslash{}to\$ Browsed at a kiosk (full screen)}.
\end{enumerate}
This is a frequent 3‑mark practical question: “Configure the presentation to run automatically for a school open day display.”
\end{propriete}

\begin{propriete}[GCE Exam Tip — Morph Transition]
The \textbf{Morph} transition (PowerPoint 2019/365) can replace dozens of manual animations. To show the growth of MTN Mobile Money agents from 2015 to 2023: duplicate the slide, resize and reposition the agent icons on the second slide, then apply Morph. PowerPoint automatically calculates the smooth movement. In written explanations, explicitly state “Morph transition” to earn the mark.
\end{propriete}

\begin{exoresolu}[GCE‑Style Practical Question]
\textbf{Task:} You have a 5‑slide presentation about “Cameroon’s Digital Economy”. 
\begin{enumerate}
\item Apply a \textbf{Fade} transition to all slides with a duration of \SI{1.00}{s} and set them to advance automatically after \SI{8}{s}.
\item On slide 3, insert the video \texttt{digital\_economy.mp4} (embedded). Set it to start \textbf{Automatically} and \textbf{Loop until Stopped}.
\item On slide 4, there is a bulleted list of four initiatives. Animate the list so each bullet appears \textbf{On Click} with a \textbf{Fade} entrance.
\item Create an Action Button on slide 1 labelled “Watch Video” that jumps directly to slide 3.
\end{enumerate}

\textbf{Solution:}
\begin{enumerate}
\item \texttt{Transitions \$\textbackslash{}to\$ Fade \$\textbackslash{}to\$ Duration: 1.00 \$\textbackslash{}to\$ Uncheck On Mouse Click \$\textbackslash{}to\$ Check After: 00:08.00 \$\textbackslash{}to\$ Apply to All.}
\item \texttt{Slide 3 \$\textbackslash{}to\$ Insert \$\textbackslash{}to\$ Video \$\textbackslash{}to\$ This Device \$\textbackslash{}to\$ digital\_economy.mp4 \$\textbackslash{}to\$ Insert. Playback \$\textbackslash{}to\$ Start: Automatically \$\textbackslash{}to\$ Loop until Stopped.}
\item \texttt{Slide 4 \$\textbackslash{}to\$ Select list \$\textbackslash{}to\$ Animations \$\textbackslash{}to\$ Fade (Entrance). Animation Pane \$\textbackslash{}to\$ Effect Options \$\textbackslash{}to\$ Text Animation \$\textbackslash{}to\$ Group text: By 1st Level Paragraphs. Verify each bullet shows Start: On Click.}
\item \texttt{Slide 1 \$\textbackslash{}to\$ Insert \$\textbackslash{}to\$ Shapes \$\textbackslash{}to\$ Action Buttons: Custom (blank) \$\textbackslash{}to\$ Draw button \$\textbackslash{}to\$ Hyperlink to: Slide... \$\textbackslash{}to\$ Slide 3 \$\textbackslash{}to\$ OK. Type “Watch Video” inside the button.}
\end{enumerate}
\end{exoresolu}

\begin{exercice}[Practice Exercise]
\begin{enumerate}
\item Open a new presentation. Create a title slide “ICT in Cameroon” and three content slides: “Mobile Money”, “E‑Government”, “Cybersecurity”.
\item Apply a \textbf{Push} transition (From Left) to all slides, duration \SI{0.75}{s}, advance \textbf{On Mouse Click}.
\item On the “Mobile Money” slide, insert an audio file \texttt{mtm\_jingle.mp3} (embedded). Set it to start \textbf{Automatically} and \textbf{Hide During Show}.
\item On the “E‑Government” slide, add a clustered bar chart comparing 2021 vs 2022 online service usage. Animate the chart \textbf{By Series} with \textbf{Wipe} entrance.
\item On the “Cybersecurity” slide, create a bulleted list of three ANTIC recommendations. Animate the list progressively \textbf{On Click} using \textbf{Appear}.
\item Add a “Home” Action Button on each content slide linking back to the title slide.
\item Save, run the show (F5), and verify all navigation and animations work as intended.
\end{enumerate}
\end{exercice}

\begin{corrige}[Exercise]
\begin{enumerate}
\item Standard slide creation.
\item \texttt{Transitions \$\textbackslash{}to\$ Push \$\textbackslash{}to\$ Effect Options: From Left \$\textbackslash{}to\$ Duration 0.75 \$\textbackslash{}to\$ Apply to All.}
\item \texttt{Insert \$\textbackslash{}to\$ Audio \$\textbackslash{}to\$ This Device \$\textbackslash{}to\$ mtm\_jingle.mp3 \$\textbackslash{}to\$ Playback \$\textbackslash{}to\$ Start: Automatically \$\textbackslash{}to\$ Hide During Show.}
\item \texttt{Insert \$\textbackslash{}to\$ Chart \$\textbackslash{}to\$ Clustered Bar \$\textbackslash{}to\$ Edit data. Animations \$\textbackslash{}to\$ Wipe \$\textbackslash{}to\$ Effect Options \$\textbackslash{}to\$ Chart Animation: By Series.}
\item \texttt{Select list \$\textbackslash{}to\$ Animations \$\textbackslash{}to\$ Appear \$\textbackslash{}to\$ Animation Pane \$\textbackslash{}to\$ Effect Options \$\textbackslash{}to\$ By 1st Level Paragraphs.}
\item \texttt{Insert \$\textbackslash{}to\$ Shapes \$\textbackslash{}to\$ Action Buttons: Home \$\textbackslash{}to\$ Draw \$\textbackslash{}to\$ Hyperlink to: Slide 1 (Title). Copy to other slides.}
\end{enumerate}
\end{corrige}

\begin{savaistu}[Did You Know? — Morph Transition in Cameroonian Classrooms]
The \textbf{Morph} transition (PowerPoint 2019/365) can replace dozens of manual animations. For example, to show the growth of MTN Mobile Money agents from 2015 to 2023, duplicate the slide, resize and reposition the agent icons on the second slide, then apply Morph. PowerPoint automatically calculates the smooth movement. This technique is a favourite in GCE practicals because it produces professional results with minimal effort — but remember to mention “Morph transition” in your written explanation to earn the mark.
\end{savaistu}

With multimedia, transitions, animations, and hyperlinks mastered, you now have the complete toolkit to produce presentations that are not only attractive but also interactive, accessible, and examination‑ready. The next section will cover delivering the presentation, printing handouts, and evaluating effectiveness — the final steps before you face the GCE board.

\coursec{Section 4: Delivering, Printing and Evaluating the Presentation}

In Section 3, you learnt how to enrich your slides with multimedia objects, transitions and animations. But a beautiful presentation that is badly delivered, badly printed or never rehearsed will fail its purpose. Imagine a candidate at the GCE practical examination who has prepared excellent slides about mobile money in Cameroon, but who does not know how to launch the slideshow, turns his back to the examiners and reads every word on the screen: all his technical work is wasted. In this final section, we learn how to \cle{run} the slideshow professionally, how to \cle{print} and \cle{export} the presentation in the right format, and how to \cle{evaluate} and \cle{deliver} it like a confident speaker.

\courssub{Running the Slideshow}

\begin{definition}[Slideshow]
A \cle{slideshow} is the full-screen playback mode of a presentation, in which slides are displayed one after another to the audience, with their transitions, animations and multimedia content active.
\end{definition}

\textbf{Starting the show.}\par
There are two essential ways to start a slideshow in Microsoft PowerPoint (and their equivalents exist in LibreOffice Impress):

\begin{itemize}
\item Press \cle{F5} (or \emph{Slide Show} $\rightarrow$ \emph{From Beginning}): the show starts from the \textbf{first slide}. This is the normal way to present to an audience.
\item Press \cle{Shift+F5} (or \emph{Slide Show} $\rightarrow$ \emph{From Current Slide}): the show starts from the slide currently selected. This is very useful when you are testing a particular slide or resuming after an interruption.
\end{itemize}

\textbf{Navigating during the show.}\par
Once the show is running, you control it with the keyboard, the mouse or discreet on-screen tools:

\begin{center}
\begin{tabularx}{\linewidth}{|l|X|}
\hline
\rowcolor{popblueL} \textbf{Action} & \textbf{How to do it} \\
\hline
Next slide / next animation & Enter, Space, Right arrow, Page Down, or left mouse click \\
\hline
Previous slide & Backspace, Left arrow, or Page Up \\
\hline
Jump to a specific slide & Type the slide number then Enter \\
\hline
End the show & \cle{Esc} \\
\hline
On-screen tools & Move the mouse to the bottom-left corner: semi-transparent buttons appear (arrows, pen, menu) \\
\hline
\end{tabularx}
\end{center}

\begin{attention}[Trap: the mouse click advances everything]
A single left click triggers the \emph{next animation}, not necessarily the next slide. If your slide has five animated bullet points, you must click five times before moving on. During a defence, candidates often panic, click too fast and skip content. Rehearse with the real file so that you know exactly how many clicks each slide requires.
\end{attention}

\courssub{Presenter Tools}

Modern presentation software offers tools designed for the person speaking, not only for the audience.

\textbf{Presenter view.}\par
When the computer is connected to a projector, \cle{Presenter View} shows two different screens: the audience sees only the current slide, while the speaker sees the current slide, the \emph{next} slide, the speaker notes and a timer. This is the professional way to present: you always know what is coming and you never lose your timing.

\textbf{Pen and highlighter during the show.}\par
During the show, press \cle{Ctrl+P} to activate the pen (or use the on-screen pen button). You can then circle a figure, underline a keyword or draw an arrow directly on the slide, exactly as a teacher does on the board. The highlighter emphasises text with a transparent colour. When you end the show, the software asks whether you want to keep or discard these annotations.

\textbf{Black or white screen.}\par
Press \cle{B} to turn the screen black, or \cle{W} to turn it white, at any moment. Press the same key again to return to the slide. This is precious when you want the audience's full attention on \emph{you} and not on the screen, for example while answering a delicate question.

\textbf{Zooming into a slide.}\par
In recent versions, you can zoom into a region of a slide during the show (on-screen magnifier tool or dedicated zoom slides) to make small details readable, for instance the figures of a budget table.

\courssub{Rehearsing Timings and Recording the Show}

\begin{methode}[Rehearsing timings]
To prepare an automatic or perfectly timed presentation:
\begin{enumerate}
\item Open the presentation and go to the \emph{Slide Show} tab.
\item Click \cle{Rehearse Timings}. The show starts full screen and a small timer toolbar appears.
\item Advance through the show exactly as you would in front of the audience, clicking when each animation or slide should appear. The timer records the duration of each slide.
\item At the end, the software displays the total time and asks whether you want to \textbf{save the timings}. Click \emph{Yes}.
\item The saved timings now appear in the Slide Sorter view under each slide, and the show can play \textbf{automatically} without any click.
\end{enumerate}
\end{methode}

\textbf{Automatic and kiosk playback.}\par
Once timings are saved, use \emph{Set Up Slide Show} to choose the playback mode. The \cle{kiosk mode} (\emph{Browsed at a kiosk}) loops the presentation continuously and ignores most keyboard input: only Esc stops it. This is ideal for an information screen at an exhibition, a school open day or a company reception hall in Douala or Yaoundé.

\textbf{Recording the slideshow.}\par
You can also \emph{Record Slide Show} to capture your voice narration together with the timings and laser-pointer movements. The result can be replayed later or exported as a video, which is useful for distance learning or for a teacher who wants students to revise a lesson at home.

\courssub{Printing the Presentation}

Printing a presentation is not only printing slides. PowerPoint offers four \cle{print layouts}, each with a precise purpose.

\begin{definition}[Handout]
A \cle{handout} is a printed version of the presentation with several slides per page (typically 2, 3, 4, 6 or 9), given to the audience so that they can follow the talk and take notes without copying the slides.
\end{definition}

\begin{itemize}
\item \textbf{Full page slides}: one slide per page. Used for archiving or for slides that must be read in detail.
\item \textbf{Handouts}: 2, 3, 4, 6 or 9 slides per page. The 3-slides layout prints lines next to each slide for handwritten notes --- the favourite for teachers and examiners.
\item \textbf{Notes pages}: one slide per page with the speaker notes printed below it. Perfect for the speaker's own preparation.
\item \textbf{Outline view}: prints only the text (titles and bullet points), without images. Ideal for a quick, cheap review of the content and structure.
\end{itemize}

\begin{popfigure}
\begin{center}
\begin{tikzpicture}[scale=1, every node/.style={font=\small}]
% Full page slide
\draw[thick, popdark] (0,0) rectangle (2.4,3.2);
\draw[fill=popblueL, draw=popblue] (0.25,0.9) rectangle (2.15,2.7);
\node[popmuted] at (1.2,0.45) {1 slide};
\node[align=center, popdark] at (1.2,-0.5) {\textbf{Full page}\\ \textbf{slide}};
% Handout 6 per page
\begin{scope}[xshift=3.6cm]
\draw[thick, popdark] (0,0) rectangle (2.4,3.2);
\foreach \x in {0.25,1.25} {
  \foreach \y in {2.15,1.25,0.35} {
    \draw[fill=popgreenL, draw=popgreen] (\x,\y) rectangle (\x+0.9,\y+0.7);
  }
}
\node[align=center, popdark] at (1.2,-0.5) {\textbf{Handout}\\ \textbf{(6 per page)}};
\end{scope}
% Notes page
\begin{scope}[xshift=7.2cm]
\draw[thick, popdark] (0,0) rectangle (2.4,3.2);
\draw[fill=popblueL, draw=popblue] (0.45,1.9) rectangle (1.95,2.95);
\draw[popmuted] (0.35,1.45) -- (2.05,1.45);
\draw[popmuted] (0.35,1.1) -- (2.05,1.1);
\draw[popmuted] (0.35,0.75) -- (2.05,0.75);
\draw[popmuted] (0.35,0.4) -- (2.05,0.4);
\node[align=center, popdark] at (1.2,-0.5) {\textbf{Notes page}};
\end{scope}
% Outline
\begin{scope}[xshift=10.8cm]
\draw[thick, popdark] (0,0) rectangle (2.4,3.2);
\draw[popdark, very thick] (0.3,2.85) -- (1.6,2.85);
\draw[popmuted] (0.5,2.5) -- (2.05,2.5);
\draw[popmuted] (0.5,2.2) -- (2.05,2.2);
\draw[popdark, very thick] (0.3,1.8) -- (1.6,1.8);
\draw[popmuted] (0.5,1.45) -- (2.05,1.45);
\draw[popmuted] (0.5,1.15) -- (2.05,1.15);
\draw[popmuted] (0.5,0.85) -- (2.05,0.85);
\draw[popdark, very thick] (0.3,0.45) -- (1.6,0.45);
\node[align=center, popdark] at (1.2,-0.5) {\textbf{Outline}\\ \textbf{view}};
\end{scope}
\end{tikzpicture}
\end{center}
\legende{The four print layouts: full page slide, handout (6 slides per page), notes page and outline view.}
\end{popfigure}

\textbf{Print range and colour.}\par
In the Print dialog you also choose:
\begin{itemize}
\item the \cle{print range}: all slides, the current slide, or a custom range such as slides 3--7;
\item the \cle{colour mode}: colour, \emph{grayscale} or \emph{pure black and white}. Grayscale saves expensive colour ink and often prints more clearly on school printers;
\item the number of copies and collation.
\end{itemize}

\begin{exemplebox}[Choosing the right layout]
For her defence on ``The impact of mobile money in Cameroon'', Ngong prints: (a) handouts with 3 slides per page for the jury, so they can annotate; (b) notes pages for herself, to rehearse; (c) the outline for her supervisor to check the logic quickly. She prints the jury's handouts in grayscale because the school printer's colour cartridge is nearly empty.
\end{exemplebox}

\courssub{Exporting and Packaging the Presentation}

Printing is not always the best way to share a presentation. Exporting creates a file that anyone can open.

\begin{propriete}[PDF export preserves the layout]
Exporting a presentation to \cle{PDF} preserves the fonts, the layout and the placement of images and media on \emph{any} computer, even one where the presentation software is not installed. Each slide becomes one page of the PDF document.
\end{propriete}

\textbf{Other export formats.}\par
\begin{itemize}
\item \textbf{Video} (\texttt{.mp4}): the slideshow, with its timings, animations and narration, becomes a video file that plays on phones, televisions and social media.
\item \textbf{Images}: each slide can be saved as a \texttt{.png} or \texttt{.jpeg} picture, useful for posting a slide on WhatsApp or inserting it into a report.
\item \textbf{Package for CD / media packaging}: this command copies the presentation \emph{together with} its linked media files (videos, sounds) and fonts into a single folder, so that nothing is missing when the file is opened on another computer.
\end{itemize}

\begin{attention}[The missing-video trap]
A video inserted as a \emph{link} is not stored inside the \texttt{.pptx} file. If you copy only the presentation to a USB flash drive and present on another computer, the video will not play. Always \emph{embed} media or use the packaging feature, and always \textbf{test the file on the actual presentation computer} before the defence.
\end{attention}

\courssub{Evaluating a Presentation: the Quality Checklist}

Before delivering, evaluate your work against six criteria. A good presentation has: a \cle{clear objective}, a \cle{logical structure}, \cle{readable slides}, a \cle{consistent design}, \cle{moderate effects} and a \cle{respected timing}.

\begin{popfigure}
\begin{center}
\begin{tikzpicture}[scale=1.05]
% radar grid
\foreach \r in {0.7,1.4,2.1,2.8} {
  \draw[popline, thin] (90:\r) \foreach \a in {150,210,270,330,30} { -- (\a:\r) } -- cycle;
}
% axes and labels
\foreach \a/\t in {90/{Objective},150/{Structure},210/{Readability},270/{Design},330/{Effects},30/{Timing}} {
  \draw[popmuted, thin] (0,0) -- (\a:2.8);
  \node[font=\small\bfseries, popdark] at (\a:3.35) {\t};
}
% performance polygon (scores out of 5 scaled to 2.8)
\draw[very thick, popblue, fill=popblueL, fill opacity=0.55]
  (90:2.52) -- (150:2.24) -- (210:1.96) -- (270:2.38) -- (330:1.68) -- (30:2.24) -- cycle;
\foreach \a/\r in {90/2.52,150/2.24,210/1.96,270/2.38,330/1.68,30/2.24} {
  \fill[popblue] (\a:\r) circle (1.6pt);
}
\node[font=\small, popmuted, align=center] at (0,-3.9) {Each axis is scored from 0 (centre) to 5 (outer ring).};
\end{tikzpicture}
\end{center}
\legende{Radar diagram of the six quality criteria. Here the presentation is strong on objective and design but weaker on effects (too many animations) and readability.}
\end{popfigure}

\begin{exoresolu}[Evaluating and improving a presentation]
Marcel has prepared 20 slides for a 5-minute talk. His slides use four different fonts, light yellow text on a white background, a different transition on every slide, and he has never rehearsed. Evaluate his presentation with the checklist and propose corrections.
\tcblower
\textbf{Solution.} Apply the six criteria:
\begin{enumerate}
\item \emph{Clear objective}: acceptable if the topic is focused, but 20 slides in 5 minutes means about 15 seconds per slide --- impossible. \textbf{Correction}: reduce to 6--8 slides.
\item \emph{Logical structure}: to verify; impose the classic plan introduction -- development -- conclusion, one idea per slide.
\item \emph{Readable slides}: light yellow on white fails the contrast rule. \textbf{Correction}: dark text on a light background (or the reverse), font size at least 24 points.
\item \emph{Consistent design}: four fonts break consistency. \textbf{Correction}: apply one theme with at most two fonts via the Slide Master.
\item \emph{Moderate effects}: a different transition per slide distracts the audience. \textbf{Correction}: keep one simple transition (e.g. Fade) for the whole show.
\item \emph{Respected timing}: never rehearsed. \textbf{Correction}: use \emph{Rehearse Timings} and adjust until the show fits within 5 minutes.
\end{enumerate}
After these corrections, Marcel's radar diagram becomes balanced, and his defence becomes credible.
\end{exoresolu}

\courssub{Delivery Skills}

The best slides in the world cannot save a poor speaker. Delivery is a skill, and it is assessed in defences and in the GCE practical context.

\begin{attention}[Common mistake: reading the slides]
The most frequent and most penalised mistake is reading the slides word for word with your back to the audience. The slides \emph{support} the speaker; they do not \emph{replace} him. The audience can read faster than you can speak: if you simply read, you add nothing and you lose eye contact, credibility and marks.
\end{attention}

\begin{methode}[Delivering like a professional]
\begin{enumerate}
\item \textbf{Face the audience}, not the screen. Glance at the slide, then speak to the people.
\item \textbf{Do not read}: use the slides as visual support and your notes pages for details.
\item \textbf{Manage time}: rehearse, watch the presenter-view timer, and know which slide you can skip if you are late.
\item \textbf{Speak clearly and slowly}, especially at the beginning when stress is highest.
\item \textbf{Handle questions}: listen fully, thank the questioner, answer briefly; if you do not know, say so honestly and promise to find out. Use the \cle{B} key to blank the screen during a long discussion.
\end{enumerate}
\end{methode}

\begin{savaistu}[Did you know?]
In Cameroon, the same delivery skills are expected everywhere: a student defending a project at the GCE, a nurse presenting a vaccination campaign in a health district, or a start-up pitching a mobile application to investors in Bonanjo all follow the same rule --- the person convinces, the slides illustrate.
\end{savaistu}

\courssub{Integration Activity}

\begin{experience}[Producing and defending a complete presentation]
In teams of three, produce and defend a complete multimedia presentation on an authentic topic, for example: ``The role of ANTIC in securing digital services in Cameroon'', ``Mobile money and financial inclusion in our town'', or ``The responsible use of social networks by students''.
\begin{enumerate}
\item Define the objective and the plan (introduction, three parts, conclusion).
\item Build 8 to 10 slides using a Slide Master for consistency; insert at least one image, one chart and one short sound or video.
\item Apply one transition to the whole show and simple entrance animations to key points.
\item Rehearse the timings until the show fits the allowed 7 minutes; save the timings.
\item Export a PDF copy and package the media on a USB flash drive.
\item Print 3-slides-per-page handouts for the jury.
\item Defend the presentation: one member presents, another handles questions, the third manages the computer. The class evaluates the defence using the six-criteria checklist and the radar diagram.
\end{enumerate}
\end{experience}

\begin{exercice}[GCE-style practice]
\begin{enumerate}
\item Give the keyboard shortcut to start a slideshow from the beginning, and the key to end it.
\item A teacher wants to give each student a printed sheet with 6 slides and room for notes. Which print layout should he choose, and in which menu?
\item Explain two advantages of exporting a presentation to PDF before sending it by e-mail.
\item During a show, which key turns the screen black, and why is this useful when answering questions?
\end{enumerate}
\end{exercice}

\begin{corrige}[Exercise 1]
\begin{enumerate}
\item \cle{F5} starts the show from the beginning; \cle{Esc} ends it.
\item He chooses the \textbf{Handouts} layout with \textbf{6 slides per page} (or 3 for ruled note lines), in \emph{File} $\rightarrow$ \emph{Print}, under the layout settings.
\item PDF preserves fonts, layout and media placement on any computer, and the recipient does not need the presentation software installed; the file is also harder to modify accidentally.
\item The \cle{B} key turns the screen black. It removes the visual distraction so the audience concentrates on the speaker during the discussion; pressing B again restores the slide.
\end{enumerate}
\end{corrige}

\begin{aretenir}[Exam tip --- GCE practical]
Memorise the essential shortcuts: \cle{F5} (start show from the beginning), \cle{Shift+F5} (from current slide), \cle{Esc} (end show), \cle{B}/\cle{W} (black/white screen), \cle{Ctrl+P} (pen). Know how to print \textbf{handouts} with a chosen number of slides per page and how to export to \textbf{PDF}: these are among the most frequent practical exam tasks. And never forget during a defence: face the jury, speak, and let the slides support you.
\end{aretenir}

\newpage

\coursec{Integration activity}

\coursec{Delivering, Printing and Evaluating the Presentation}

\courssub{Aims of the activity}
This integration activity is designed to consolidate all competencies acquired in Lessons 65, 66 and 67. By the end of this activity, you must be able to:
\begin{itemize}
    \item \textbf{Plan} a presentation project: analyse the audience, define a clear communication objective, structure a logical outline and validate it with a storyboard.
    \item \textbf{Implement} a consistent visual identity using the \cle{Slide Master}: apply a custom theme, insert a persistent logo, configure footers with automatic slide numbering and date, and define master layouts for title, content and section header slides.
    \item \textbf{Produce} professional slides: manipulate text boxes, bulleted lists, \cle{SmartArt} graphics, tables and charts (linked to a data source) while respecting readability rules (contrast, font size, alignment, 6$\times$6 rule).
    \item \textbf{Integrate multimedia} elements purposefully: insert and trim a local video clip or audio narration, set playback options (automatic/on click, loop, hide while not playing) and compress media to manage file size.
    \item \textbf{Apply transitions and animations} with moderation: choose a single, subtle transition scheme for the whole deck; apply progressive entrance animations to lists and objects; use the \cle{Animation Pane} to sequence effects (Start: On Click / After Previous) and adjust timings.
    \item \textbf{Prepare delivery}: rehearse with the \cle{Rehearse Timings} tool, create speaker notes, configure \cle{Presenter View} for dual-screen support, and set up a kiosk or speaker-led show type.
    \item \textbf{Export and print} professional outputs: print 3-slides-per-page handouts with note lines, export a high-fidelity \cle{PDF} for distribution, and package the presentation for CD/USB (\cle{Package Presentation for CD}).
    \item \textbf{Evaluate} a presentation critically using a standardised checklist covering structure, visual design, content accuracy, technical execution and oral delivery.
\end{itemize}

\courssub{Situation}
\textbf{Context:} The \emph{Government Bilingual High School (GBHS) Mbankomo} has been invited to showcase its achievements at the \emph{National Education Open Days (Journées Portes Ouvertes de l'Éducation)} organised by MINEDUB at the \emph{Palais des Congrès} in \textbf{Yaoundé}, next month. The Principal, Mr. \textsc{Fonkou}, will man the school's stand and meet prospective parents, sponsors (local businesses, NGOs) and educational partners.

\textbf{Problem:} The Principal needs a \textbf{10-slide multimedia presentation} that runs automatically on a loop at the stand (kiosk mode) but can also be delivered live in 5 minutes during a scheduled slot on the main stage. The presentation must convince the audience that GBHS Mbankomo is the school of choice for academic excellence, holistic development and modern facilities.

\textbf{Constraints and Requirements:}
\begin{enumerate}
    \item \textbf{Teamwork:} Work in groups of \textbf{three students} (roles: Project Manager/Designer, Content Researcher/Data Analyst, Multimedia Technician/Presenter).
    \item \textbf{Audience:} Parents (Francophone/Anglophone), Sponsors (MTN Foundation, Orange Cameroon, local entrepreneurs), MINEDUB inspectors.
    \item \textbf{Objective:} \emph{Inform} about results and facilities, \emph{Persuade} to enrol/partner, \emph{Reassure} on discipline and values.
    \item \textbf{Length:} Exactly \textbf{10 slides} (including title and closing slides). Live delivery: \textbf{5 minutes} maximum.
    \item \textbf{Mandatory Content Elements:}
        \begin{itemize}
            \item At least \textbf{one original photo} of the school (building, lab, library, or sports field).
            \item One \textbf{Clustered Column Chart} showing GCE Advanced Level pass rates (\% pass) for the last three sessions (2022, 2023, 2024) for both \emph{Arts} and \emph{Science} series.
            \item One \textbf{Table} summarising the \textbf{Annual School Fees Structure} (Tuition, PTA, Uniform, Lab/Computer fees) for Form 1 to Upper Sixth.
            \item One \textbf{multimedia clip}: either a 30-second \textbf{video} (school choir, lab experiment, sports action) or an \textbf{audio narration} of the Principal's welcome message (recorded via smartphone/PC microphone).
        \end{itemize}
    \item \textbf{Technical Specifications:}
        \begin{itemize}
            \item \textbf{Slide Master:} Custom theme (school colours: Navy Blue \#002B5C and Gold \#C5962C), School Logo top-right, Footer with "GBHS Mbankomo | Yaoundé Open Days 2025 | Slide \#" and automatic date.
            \item \textbf{Transitions:} One uniform transition applied to \emph{All Slides} (e.g., \cle{Fade} or \cle{Push}, Duration 1.00s). No sound on transition.
            \item \textbf{Animations:} Progressive entrance (e.g., \cle{Appear} or \cle{Wipe}) for all bulleted lists (By Paragraph: Level 1). Charts/Table: By Series/Category. Entrance only; \emph{no} Emphasis or Exit effects.
            \item \textbf{Rehearsal:} Timings recorded via \cle{Rehearse Timings}. Total show $\le$ 5:00 min.
            \item \textbf{Outputs:} (a) Print \textbf{Handouts (3 slides/page)} with lines for notes. (b) Export \textbf{PDF} (Standard quality). (c) \textbf{Package for CD/USB} (linked media included).
        \end{itemize}
    \item \textbf{Defense \& Evaluation:} Each team presents live (5 min) to the class. Peers evaluate using the \textbf{Peer Evaluation Checklist} (provided below). The teacher validates the best presentation for the Principal.
\end{enumerate}

\courssub{Tasks}
\textbf{Phase 1: Planning \& Storyboarding (Paper-based, 30 min)}
\begin{enumerate}
    \item \textbf{Audience \& Objective Analysis:} Write a 3-sentence \cle{Communication Objective} (SMART: Specific, Measurable, Achievable, Relevant, Time-bound). Identify 3 key concerns of the \emph{Parents} and 2 of the \emph{Sponsors}.
    \item \textbf{Outline \& Storyboard:} On A4 paper, sketch a 10-slide storyboard. For each slide: (a) Slide Title, (b) Layout type (Title, Title+Content, Two Content, Section Header, Blank), (c) Key message (max 6 words), (d) Visual/Multimedia element planned, (e) Speaker notes prompt. Get teacher \textbf{approval signature} before touching the computer.
\end{enumerate}

\textbf{Phase 2: Slide Master \& Theme Construction (20 min)}
\begin{enumerate}[resume]
    \item Open a blank presentation. Go to \texttt{View > Slide Master}.
    \item \textbf{Theme Colours:} \texttt{Slide Master > Background > Colours > Customize Colours}. Set: Text/Background Dark 1 = Navy \#002B5C; Accent 1 = Gold \#C5962C; Accent 2 = White \#FFFFFF; Hyperlink = Gold. Save as \texttt{GBHS\_Theme}.
    \item \textbf{Fonts:} \texttt{Fonts > Customize Fonts}. Heading: \texttt{Calibri Light (Headings)}; Body: \texttt{Calibri (Body)}.
    \item \textbf{Master Layout (Top Master):} Insert school logo (top-right, height 1.2 cm). Insert \texttt{Footer} placeholder: Left: \texttt{GBHS Mbankomo | Yaoundé Open Days 2025}; Center: \texttt{Date (Fixed: 15 March 2025)}; Right: \texttt{Slide \#}. Apply \texttt{Gold} font colour to footer text. Untick \texttt{Title Slide} footer if desired.
    \item \textbf{Custom Layouts:} Duplicate \texttt{Title and Content} layout $\rightarrow$ Rename \texttt{Chart\_Layout}. Add a placeholder \texttt{Content > Chart} on right, text on left. Duplicate for \texttt{Table\_Layout}. Close Master View.
\end{enumerate}

\textbf{Phase 3: Slide Development \& Content Population (60 min)}
\begin{enumerate}[resume]
    \item \textbf{Slide 1 (Title Layout):} "GBHS Mbankomo: Excellence in Bilingual Education". Subtitle: "Yaoundé Open Days 2025 | Preparing Global Citizens". Background: Full-bleed school gate photo (send to back).
    \item \textbf{Slide 2 (Section Header):} "Our Vision \& Mission". Body: 3 bullets (Vision, Mission, Core Values).
    \item \textbf{Slide 3 (Title + Content + Photo):} "Modern Facilities for 21st Century Learning". Left: 4 bullets (Science Labs, ICT Centre, Library, Sports Complex). Right: Insert school photo (lab or library). Apply \cle{Picture Format > Corrections > Sharpen/Soften} if needed.
    \item \textbf{Slide 4 (Chart\_Layout):} "GCE A/L Pass Rates: 3-Year Trend (2022--2024)". Insert Clustered Column Chart. \textbf{Data:} See Model Answer. Chart Title: "Pass Rate \%". Legend: Bottom. Data Labels: Outside End. Style: Style 8 (Gold columns on Navy).
    \item \textbf{Slide 5 (Table\_Layout):} "Affordable Fee Structure 2024/2025". Insert Table (6 rows $\times$ 5 cols). Header Row: Class, Tuition (FCFA), PTA (FCFA), Uniform (FCFA), Total (FCFA). Fill realistic data (see Model). Apply \texttt{Table Style Medium 2} (Gold header). Set \texttt{Layout > AutoFit > Fixed Column Width}.
    \item \textbf{Slide 6 (Two Content):} "Co-Curricular Excellence". Left: SmartArt \texttt{Vertical Box List} (Clubs: Robotics, Debate, Environment, Journalism). Right: Photo collage (3 pics grouped) or icons.
    \item \textbf{Slide 7 (Title + Content):} "Strategic Partnerships". Logos of MTN Foundation, Orange Cameroon, MINEDUB, Local Council (Insert > Pictures > This Device). Align horizontally, distribute horizontally.
    \item \textbf{Slide 8 (Blank Layout + Multimedia):} "Hear from Our Community". Insert \textbf{Video} (from file: \texttt{choir\_performance.mp4}) OR \textbf{Audio} (from file: \texttt{principal\_welcome.m4a}). \textbf{Playback:} Start \texttt{Automatically} (kiosk) / \texttt{On Click} (live). Check \texttt{Loop until Stopped} (kiosk only). \texttt{Video Format > Poster Frame} set to a nice frame.
    \item \textbf{Slide 9 (Title + Content):} "Admissions Open for 2025/2026". Bullets: Entry classes, Requirements, Deadline, Contact (Phone/WhatsApp: +237 6XX XX XX XX, Email: admissions@gbhsmbankomo.cm). Icon for phone/email.
    \item \textbf{Slide 10 (Title Only / Closing):} "Thank You / Merci / Mèsi". School Logo large centre. "Visit our Stand \#B12". QR Code image linking to school website (generated online).
\end{enumerate}

\textbf{Phase 4: Transitions, Animations \& Timings (30 min)}
\begin{enumerate}[resume]
    \item \textbf{Transitions:} \texttt{Transitions > Fade} (or Push). \texttt{Effect Options: Through Black} (if Fade) / \texttt{From Right} (if Push). Duration: \texttt{1.00 s}. \texttt{Apply To All}. Uncheck \texttt{On Mouse Click} (for kiosk); Check \texttt{After: 00:00.00} (will be overwritten by rehearsal).
    \item \textbf{Animations:} Select Slide 2 content placeholder $\rightarrow$ \texttt{Animations > Appear}. \texttt{Effect Options > By Paragraph}. Open \texttt{Animation Pane}. Select all list animations $\rightarrow$ \texttt{Start: After Previous}. Delay: \texttt{0.20 s}. Repeat for Slides 3, 4 (Chart: By Series), 5 (Table: By Row), 6, 7, 9. \textbf{Remove} all animations from Slide 1 and 10.
    \item \textbf{Rehearse Timings:} \texttt{Slide Show > Rehearse Timings}. Deliver script naturally. Pause on Slide 8 for video/audio to play fully. Save timings. Verify total $\le$ 5:00 min in \texttt{Slide Sorter} view.
    \item \textbf{Set Up Show:} \texttt{Slide Show > Set Up Slide Show}. Choose \texttt{Browsed at a kiosk (full screen)} for stand version; \texttt{Presented by a speaker (full screen)} for live version. Save two copies if needed: \texttt{GBHS\_Kiosk.pptx} and \texttt{GBHS\_Live.pptx}.
\end{enumerate}

\textbf{Phase 5: Outputs, Packaging \& Defense (20 min)}
\begin{enumerate}[resume]
    \item \textbf{Handouts:} \texttt{File > Print > Settings > Full Page Slides > Handouts > 3 Slides}. Check \texttt{Frame Slides}, \texttt{Print Comments} (if any), \texttt{Print Hidden Slides} (off). Print to PDF or Paper.
    \item \textbf{Export PDF:} \texttt{File > Export > Create PDF/XPS Document > Create PDF/XPS}. Options: \texttt{Standard (publishing online and printing)}. Publish.
    \item \textbf{Package for CD:} \texttt{File > Export > Package Presentation for CD > Package for CD}. Name: \texttt{GBHS\_OpenDays\_Package}. Ensure \texttt{Linked files} included. Copy to USB.
    \item \textbf{Live Defense:} Team presents (5 min). Class evaluates using Checklist. Teacher moderates.
\end{enumerate}

\courssub{Model answer}
\textbf{1. Planning Artefacts (Sample)}

\begin{center}
\begin{tabularx}{\linewidth}{|l|X|}\hline
\rowcolor{popblueL} \textbf{Element} & \textbf{Description} \\ \hline
\textbf{Communication Objective} & \emph{By the end of the 5-minute presentation, 80\% of visiting parents at the Yaoundé Open Days will be able to cite three distinct academic or infrastructural strengths of GBHS Mbankomo and request an admission form or scan the QR code for more information.} \\ \hline
\textbf{Parent Concerns} & 1. Academic results (GCE pass rates). 2. Discipline \& moral values. 3. Cost/Value for money (fees vs facilities). \\ \hline
\textbf{Sponsor Concerns} & 1. Visibility/ROI (branding on stand, uniforms). 2. Impact metrics (students reached, STEM labs equipped). \\ \hline
\end{tabularx}
\end{center}

\textbf{Storyboard Outline (Approved by Teacher: \underline{\hspace{3cm}} Date: \underline{\hspace{2cm}})}

\begin{center}
\begin{tabularx}{\linewidth}{|c|l|l|X|X|l|}\hline
\rowcolor{popblueL} \textbf{\#} & \textbf{Title} & \textbf{Layout} & \textbf{Key Message} & \textbf{Visual/Media} & \textbf{Notes Prompt} \\ \hline
1 & GBHS Mbankomo: Excellence... & Title & Brand identity & School Gate Photo (Full bleed) & "Welcome, I am [Name]..." \\ \hline
2 & Our Vision \& Mission & Section Header & Guiding principles & None (Clean) & "Three pillars..." \\ \hline
3 & Modern Facilities & Title+Content+Pic & Conducive environment & Lab Photo (Right) & "Invested 50M FCFA..." \\ \hline
4 & GCE Results 2022-2024 & Chart\_Layout & Consistent excellence & Clustered Column Chart & "Science 92\%, Arts 88\%..." \\ \hline
5 & Fee Structure & Table\_Layout & Transparent, affordable & Fees Table & "All-inclusive, no hidden costs" \\ \hline
6 & Co-Curricular Life & Two Content & Holistic development & SmartArt + Icons & "Robotics won Regionals..." \\ \hline
7 & Partnerships & Title+Content & Trusted by leaders & Partner Logos & "MTN equipped ICT Lab..." \\ \hline
8 & Community Voices & Blank + Media & Authentic testimony & Video/Audio Clip & "Play video... pause for effect" \\ \hline
9 & Admissions 2025/26 & Title+Content & Call to Action & Icons (Phone, Email, QR) & "Forms available at Stand B12" \\ \hline
10 & Thank You & Title Only & Lasting impression & Logo + QR Code & "Questions? Visit Stand B12" \\ \hline
\end{tabularx}
\end{center}

\textbf{2. Slide Master Configuration Details}
\begin{itemize}
    \item \textbf{Colour Palette (RGB):} Navy \texttt{0, 43, 92} (\#002B5C); Gold \texttt{197, 150, 44} (\#C5962C); White \texttt{255, 255, 255}; Light Grey \texttt{230, 230, 230} (for table alt rows).
    \item \textbf{Footer Text Boxes (Master):} Left: \texttt{GBHS Mbankomo | Yaoundé Open Days 2025} (Font: Calibri 10 pt, Gold, Left Align). Center: \texttt{15 March 2025} (Fixed date). Right: \texttt{Slide \#} (Slide Number placeholder, Gold).
    \item \textbf{Logo Placement:} Top Master $\rightarrow$ \texttt{Insert > Pictures > Logo.png} $\rightarrow$ Position: Horizontal Alignment: Right, Vertical: Top. Margin: 0.5 cm. \texttt{Lock Aspect Ratio}.
\end{itemize}

\textbf{3. Slide 4 -- Chart Data (Source: School Archives, verified by Principal)}
\begin{center}
\begin{tabular}{|c|c|c|}\hline
\rowcolor{popblueL} \textbf{Session} & \textbf{Arts Series (\% Pass)} & \textbf{Science Series (\% Pass)} \\ \hline
2022 & 78 & 85 \\ \hline
2023 & 82 & 89 \\ \hline
2024 & 88 & 92 \\ \hline
\end{tabular}
\end{center}
\textit{Chart Construction:} Insert $\rightarrow$ Chart $\rightarrow$ Clustered Column. Paste data above. \textbf{Chart Elements:} Axes (Primary Vertical: 0--100\%, Major Gridlines off), Chart Title: "GCE A/L Pass Rate Trend (2022--2024)", Legend: Bottom. \textbf{Format Data Series:} Arts $\rightarrow$ Fill: Gold (\#C5962C); Science $\rightarrow$ Fill: Navy (\#002B5C). Gap Width: 100\%. Data Labels: Value, Outside End, Font: Calibri 10 pt, Bold, Navy.

\textbf{4. Slide 5 -- Fee Structure Table (Sample Data FCFA)}
\begin{center}
\begin{tabular}{|l|r|r|r|r|}\hline
\rowcolor{popgold} \textbf{Class} & \textbf{Tuition} & \textbf{PTA Levy} & \textbf{Uniform (Set)} & \textbf{Total Annual} \\ \hline
Form 1 -- 2 & 120,000 & 15,000 & 25,000 & 160,000 \\ \hline
Form 3 -- 4 & 130,000 & 15,000 & 25,000 & 170,000 \\ \hline
Form 5 & 140,000 & 15,000 & 25,000 & 180,000 \\ \hline
Lower Sixth (Arts) & 150,000 & 20,000 & 30,000 & 200,000 \\ \hline
Lower Sixth (Science) & 160,000 & 20,000 & 30,000 & 210,000 \\ \hline
Upper Sixth (Arts) & 155,000 & 20,000 & 0* & 175,000 \\ \hline
Upper Sixth (Science) & 165,000 & 20,000 & 0* & 185,000 \\ \hline
\end{tabular}
\end{center}
\emph{* Uniform purchased in Lower Sixth.} \textbf{Table Style:} Medium 2 (Gold Header, Banded Rows). \textbf{Animation:} \texttt{Animations > Wipe > Effect Options: From Left > By Row}. Start: After Previous, Delay 0.15s.

\textbf{5. Slide 8 -- Multimedia Configuration}
\begin{itemize}
    \item \textbf{Video File:} \texttt{choir\_performance.mp4} (H.264/MP4, 720p, 28 sec, 12 MB). Inserted via \texttt{Insert > Media > Video > This Device}.
    \item \textbf{Playback Tab Settings (Kiosk Version):} \texttt{Start: Automatically}; \texttt{Play Full Screen: Checked}; \texttt{Hide While Not Playing: Checked}; \texttt{Loop until Stopped: Checked}; \texttt{Rewind after Playing: Checked}.
    \item \textbf{Playback Tab Settings (Live Version):} \texttt{Start: On Click}; \texttt{Loop until Stopped: Unchecked}.
    \item \textbf{Poster Frame:} Scrub to 00:05 (choir smiling) $\rightarrow$ \texttt{Video Format > Poster Frame > Current Frame}.
    \item \textbf{Compression:} \texttt{File > Info > Media Size and Performance > Compress Media > HD (720p)}. Target final .pptx < 50 MB.
\end{itemize}

\textbf{6. Animation Pane Configuration (Representative Sample for Slide 4 Chart)}
\begin{center}
\begin{tabular}{|l|l|l|l|l|}\hline
\rowcolor{popblueL} \textbf{Order} & \textbf{Object} & \textbf{Effect} & \textbf{Start} & \textbf{Duration / Delay} \\ \hline
1 & Chart Title & Appear & After Previous & 0.50s / 0.00s \\ \hline
2 & Chart Area (Background) & Fade & After Previous & 0.50s / 0.00s \\ \hline
3 & Series "Arts" (3 columns) & Wipe (From Bottom) & After Previous & 0.50s / 0.20s \\ \hline
4 & Series "Science" (3 columns) & Wipe (From Bottom) & After Previous & 0.50s / 0.20s \\ \hline
5 & Legend & Fade & After Previous & 0.30s / 0.00s \\ \hline
\end{tabular}
\end{center}
\textit{Note:} All text placeholders on content slides use \texttt{Appear > By Paragraph (Level 1) > Start: After Previous > Delay: 0.20s}. No animations on Slide 1 (Title) and Slide 10 (Closing).

\textbf{7. Rehearsal Timings Log (Target: 5:00 min = 300 sec)}
\begin{center}
\begin{tabular}{|c|l|r|r|}\hline
\rowcolor{popblueL} \textbf{Slide} & \textbf{Content} & \textbf{Elapsed (mm:ss)} & \textbf{Notes} \\ \hline
1 & Title & 0:00 -- 0:15 & Walk to centre, smile \\ \hline
2 & Vision/Mission & 0:15 -- 0:45 & 3 bullets, 10s each \\ \hline
3 & Facilities + Photo & 0:45 -- 1:20 & Point to photo \\ \hline
4 & Chart Results & 1:20 -- 2:10 & Highlight 2024 peak \\ \hline
5 & Fees Table & 2:10 -- 2:50 & Row-by-row reveal \\ \hline
6 & Co-Curricular & 2:50 -- 3:20 & Mention Robotics win \\ \hline
7 & Partnerships & 3:20 -- 3:45 & Name MTN, Orange \\ \hline
8 & Video/Audio & 3:45 -- 4:15 & \textbf{Let media play fully (30s)} \\ \hline
9 & Admissions & 4:15 -- 4:45 & Point to QR code \\ \hline
10 & Thank You & 4:45 -- 5:00 & "Questions? Stand B12" \\ \hline
\end{tabular}
\end{center}
\textbf{Action:} If rehearsal exceeds 5:00, reduce speaker notes on Slides 3, 6, 7; increase animation speed (Duration 0.30s).

\textbf{8. Output Generation Checklist}
\begin{itemize}
    \item \textbf{Done:} \textbf{Handouts (3 slides/page):} \texttt{File > Print > Printer: Microsoft Print to PDF > Settings: 3 Slides Handouts > Frame Slides: On > Save as \textbackslash{}texttt\{GBHS\_Handouts.pdf\}}. Verify: 4 pages (Slides 1-3, 4-6, 7-9, 10+blank), note lines present.
    \item \textbf{Done:} \textbf{Full PDF:} \texttt{File > Export > Create PDF/XPS > Options: Publish what: Slides, Include: Document properties, Document structure tags for accessibility > Publish as \textbackslash{}texttt\{GBHS\_Presentation.pdf\}}.
    \item \textbf{Done:} \textbf{Package for CD:} \texttt{File > Export > Package Presentation for CD > Name: GBHS\_OpenDays\_Package > Options: Linked files (checked), Embed TrueType fonts (checked) > Copy to Folder/USB}. Verify folder contains \texttt{GBHS\_Live.pptx}, \texttt{choir\_performance.mp4}, \texttt{principal\_welcome.m4a}, \texttt{Logo.png}.
\end{itemize}

\textbf{9. Peer Evaluation Checklist (Used during Defense)}
Each peer group scores the presenting team 1 (Poor) to 5 (Excellent).

\begin{center}
\begin{tabularx}{\linewidth}{|l|X|ccccc|}\hline
\rowcolor{popblueL} \textbf{Criteria} & \textbf{Descriptors} & \textbf{1} & \textbf{2} & \textbf{3} & \textbf{4} & \textbf{5} \\ \hline
\textbf{Structure \& Logic} & Clear objective, logical flow, strong opening/closing, signposting. & & & & & \\ \hline
\textbf{Visual Design} & Consistent Master (logo, footer, colours), readable fonts ($\ge$24pt), high contrast, 6$\times$6 rule respected, quality images. & & & & & \\ \hline
\textbf{Content Accuracy} & Correct data (chart/table), relevant info, no spelling/grammar errors, citations if needed. & & & & & \\ \hline
\textbf{Multimedia Integration} & Video/Audio plays smoothly, relevant, trimmed, formatted (poster frame), file size managed. & & & & & \\ \hline
\textbf{Transitions/Animations} & Uniform transition, subtle; animations progressive (by paragraph/series), purposeful, not distracting, timed well. & & & & & \\ \hline
\textbf{Timing \& Pacing} & Total $\le$ 5:00 min, slides advance smoothly, speaker not rushed, media accommodated. & & & & & \\ \hline
\textbf{Oral Delivery} & Eye contact, voice projection, bilingual (En/Fr) elements, handles Q\&A, uses Presenter View/notes naturally. & & & & & \\ \hline
\textbf{Outputs Ready} & Handouts (3-up), PDF, Package folder complete and functional on USB. & & & & & \\ \hline
\end{tabularx}
\end{center}

\textbf{Scoring:} Total / 40. \textbf{Pass mark: 28/40}. The team with the highest score (teacher moderates) submits the final \texttt{GBHS\_Kiosk.pptx} and \texttt{GBHS\_Live.pptx} to the Principal via the school's administrative email (\texttt{principal@gbhsmbankomo.cm}) and USB key.

\begin{attention}[Common Pitfalls in Integration Activities]
\begin{itemize}
\item \textbf{Forgetting the Master:} Formatting each slide individually leads to inconsistency (misaligned logos, varying fonts). \emph{Always} use Slide Master first.
\item \textbf{Linking vs Embedding Media:} Inserting a video from a local path creates a \emph{link} by default in older versions; in modern 365 it embeds. \textbf{Always verify} in \texttt{File > Info > Optimize Media Compatibility} and use \textbf{Package for CD} to ensure files travel with the .pptx.
\item \textbf{Animation Overload:} Using "Fly In", "Bounce", "Spin" on every bullet distracts the jury. \emph{Stick to one entrance effect (Appear/Wipe/Fade) for the whole deck.}
\item \textbf{Chart Junk:} 3D charts, excessive gridlines, legend on right taking space. \emph{Use 2D Clustered Column, remove vertical gridlines, place legend at bottom.}
\item \textbf{Ignoring Kiosk vs Live:} A kiosk show \emph{must} have \texttt{Transitions > After: [time]} and \texttt{Video > Start: Automatically}. A live show \emph{must} have \texttt{On Mouse Click} and \texttt{Video > Start: On Click}. Save two versions.
\item \textbf{Font Substitution:} Using a fancy downloaded font (e.g., "African Sunset") that is not on the presentation PC at the Palais des Congrès. \emph{Use system fonts (Calibri, Arial, Segoe UI) or Embed Fonts (File > Options > Save > Embed fonts in the file).}
\end{itemize}
\end{attention}

\begin{savaistu}[Cameroon Context: Digital Skills for Employability]
The ability to produce a polished, self-running multimedia presentation is a core \cle{21st-century skill} highlighted in Cameroon's \emph{National ICT Policy} and the \emph{GCE A-Level ICT Syllabus}. In the professional world — whether pitching a startup to the \emph{Yaoundé Business Plan Competition}, presenting a project report at \emph{CAMTEL} or \emph{ANTIC}, or defending a thesis at \emph{University of Yaoundé I} — the standards are identical: \textbf{structure, visual consistency, technical reliability, and time management}. This activity simulates a real \emph{Request for Proposal (RFP)} scenario where the "client" (the Principal) has strict specifications. Mastering the \cle{Slide Master} separates amateurs (who format manually) from professionals (who design systems). The \cle{Package for CD} feature ensures your work survives the "it works on my laptop" syndrome — a classic embarrassment at national events like the \emph{Foire Internationale de l'Entreprise (FIE)} or \emph{Journées de l'Excellence}. Treat this exercise as a portfolio piece for your future career.
\end{savaistu}

\newpage

\coursec{Methods and worked examples}

\coursec{Getting Started with Presentation Software and Producing Slides}

A \cle{presentation software} (Microsoft PowerPoint, LibreOffice Impress, Google Slides) is used to produce a sequence of \cle{slides}: pages displayed on a screen or projected, combining text, images, charts, sound and video to support an oral talk. At the GCE Advanced Level you must be able to \emph{produce}, \emph{illustrate}, \emph{animate} and \emph{deliver} a slide presentation, and to \emph{evaluate} one critically.

\begin{definition}[Slide, layout, presentation]
A \cle{slide} is a single page of a presentation. A \cle{layout} is a predefined arrangement of \emph{placeholders} (title, text, picture, chart\ldots) on a slide, e.g. \emph{Title Slide}, \emph{Title and Content}, \emph{Two Content}, \emph{Comparison}, \emph{Blank}. A \cle{presentation} is the whole file: an ordered set of slides plus its design (theme, master), its effects (transitions, animations) and its media.
\end{definition}

\courssub{Skill 1: Creating a presentation and choosing an appropriate layout}

\begin{methode}[Producing a slide presentation from scratch]
To produce a clean slide presentation with software such as Microsoft PowerPoint or LibreOffice Impress:
\begin{enumerate}
\item \textbf{Plan before you click.} On paper, list the title, the sections and one idea per slide. A good rule is \emph{one slide = one message}.
\item \textbf{Launch the software} and choose \emph{New presentation} (blank or from a template/design theme).
\item \textbf{Choose the slide layout} (Title Slide, Title and Content, Two Content, Comparison, Blank\ldots) from the \emph{Layout} menu \emph{before} typing; changing it later can displace your content.
\item \textbf{Enter the text} in the placeholders. Keep phrases short (keywords, not paragraphs).
\item \textbf{Insert new slides} with \emph{New Slide} (shortcut \texttt{Ctrl+M}) and duplicate or delete slides from the slide pane.
\item \textbf{Save early and often} (\texttt{Ctrl+S}) in the correct format: \texttt{.pptx} (PowerPoint) or \texttt{.odp} (Impress), and keep a copy on a USB key.
\item \textbf{Check the result} in Slide Sorter view, then run the show with \texttt{F5}.
\end{enumerate}
\textbf{Reflexes:} title on every slide; maximum about 6 lines per slide and 6 words per line (the ``6$\times$6 rule''); save under a meaningful file name, e.g. \texttt{Water\_Project\_Group3.pptx}.
\end{methode}

\begin{attention}[Common beginner mistakes]
\begin{itemize}
\item Typing whole paragraphs on a slide: the audience reads instead of listening. Details belong to the \emph{speech} and the \emph{speaker notes}.
\item Choosing the layout \emph{after} typing: content may jump or disappear.
\item Saving only at the end: a power cut (frequent during the rainy season!) can destroy an hour of work. Save every few minutes.
\end{itemize}
\end{attention}

\begin{exoresolu}[Creating a first presentation]
\textbf{Statement (GCE style).} For the end-of-year open day, your school asks each Upper Sixth class to present one science project. Your group must prepare a 5-slide presentation entitled \emph{``Clean Water for Buea''}. (a) Give the layout you would choose for slide 1 and for slide 3 (which lists three causes of water pollution). (b) Describe the steps to add a fourth slide between slides 3 and 4. (c) State two good practices for the text on the slides.
\tcblower
\textbf{Solution.}
\begin{enumerate}
\item[(a)] \textbf{Slide 1} is the opening slide: choose the \cle{Title Slide} layout, which contains a title placeholder (\emph{Clean Water for Buea}) and a subtitle placeholder (group members, class, date). \textbf{Slide 3} presents a list of three causes: choose the \cle{Title and Content} layout, with the title \emph{Causes of water pollution} and a bulleted list in the content placeholder (domestic waste, chemical fertilisers, leaking latrines).
\item[(b)] To insert a slide between slides 3 and 4: click on slide 3 in the slide pane (left panel) or in Slide Sorter view, then click \emph{Home $\rightarrow$ New Slide} (or press \texttt{Ctrl+M}). The new slide is inserted \emph{after} the selected one and automatically becomes slide 4; the old slide 4 is renumbered 5. Then choose its layout and type the content.
\item[(c)] Good practices for text: (i) use short phrases or keywords rather than full paragraphs, respecting roughly the 6$\times$6 rule; (ii) keep a large, readable font size (at least 24 pt for body text) and the same font throughout the presentation.
\end{enumerate}
\textbf{Examiner's remark:} marks are given for naming the layouts \emph{precisely} (``Title Slide'', ``Title and Content'') and for stating that insertion happens \emph{after} the selected slide.
\end{exoresolu}

\coursec{Designing Attractive and Consistent Slides}

\courssub{Skill 2: Applying themes, masters and consistent formatting}

\begin{definition}[Theme and Slide Master]
A \cle{theme} is a coordinated set of colours, fonts and background styles applied to all slides at once. The \cle{Slide Master} is the model slide from which all slides inherit their formatting: any element placed on the master (logo, font, footer, background) is repeated automatically on every slide that uses it.
\end{definition}

\begin{methode}[Designing attractive and consistent slides]
To obtain a professional, uniform look:
\begin{enumerate}
\item Apply a \cle{theme} (Design tab): it fixes the colour palette, fonts and background of all slides at once.
\item For deeper control, open the \cle{Slide Master} (\emph{View $\rightarrow$ Slide Master}): any change made on the master (logo, font, footer, background) is repeated on every slide using that master.
\item Add \emph{slide numbers}, date and footer via \emph{Insert $\rightarrow$ Header \& Footer}.
\item Use \emph{high contrast}: dark text on light background or the reverse; avoid red/green combinations (colour-blind audience).
\item Align and space objects with \emph{Arrange $\rightarrow$ Align / Distribute}; group related objects.
\item Limit yourself to \textbf{two fonts} and \textbf{three colours} maximum.
\end{enumerate}
\textbf{Key idea:} consistency is achieved through the \emph{master}, not by formatting each slide by hand.
\end{methode}

\begin{exoresolu}[Using the Slide Master]
\textbf{Statement.} A student has typed a 20-slide presentation on \emph{HIV/AIDS sensitisation in the North West Region}. The school now requires that every slide display the school logo at the top right, slide numbers at the bottom right, and the font \emph{Calibri}. (a) Explain the most efficient way to satisfy these three requirements. (b) The student formatted slide 7 manually with a different font before using the master, and slide 7 does not change. Explain why and propose a remedy. (c) Give two reasons why consistency of design matters in a presentation.
\tcblower
\textbf{Solution.}
\begin{enumerate}
\item[(a)] Open \emph{View $\rightarrow$ Slide Master}. On the master slide: (i) insert the logo picture (\emph{Insert $\rightarrow$ Picture}) and position it at the top right; (ii) select the title and body placeholders and set the font to Calibri; (iii) close the Master view, then use \emph{Insert $\rightarrow$ Header \& Footer}, tick \emph{Slide number}, and click \emph{Apply to All}. Because all 20 slides inherit from the master, the logo, font and numbers appear everywhere in a single operation instead of 20 manual edits.
\item[(b)] Manual (local) formatting \emph{overrides} the master: slide 7 keeps its directly applied font. Remedy: select slide 7 and use \emph{Reset} (Reset Slide / Reset Layout) to remove local formatting so the slide follows the master again.
\item[(c)] Consistency (i) helps the audience focus on the \emph{message} instead of being distracted by changing styles, and (ii) gives a \emph{professional, credible} image of the presenter and the school.
\end{enumerate}
\textbf{Examiner's remark:} the phrase ``manual formatting overrides the master'' is the key marking point in (b).
\end{exoresolu}

\courssub{Skill 3: Inserting and formatting illustrative objects}

\begin{methode}[Illustrating slides with images, tables, charts and SmartArt]
\begin{enumerate}
\item \textbf{Images:} \emph{Insert $\rightarrow$ Pictures}; resize by a \emph{corner} handle to keep proportions; use \emph{Compress Pictures} to reduce file size.
\item \textbf{Tables:} \emph{Insert $\rightarrow$ Table}; use them for exact numerical data (e.g. marks, prices).
\item \textbf{Charts:} \emph{Insert $\rightarrow$ Chart}; choose the type according to the data: \emph{column/bar} for comparisons, \emph{line} for evolution over time, \emph{pie} for shares of a total (percentages).
\item \textbf{SmartArt / diagrams:} for processes, hierarchies, cycles (e.g. the water cycle).
\item Always add a \emph{title or caption} and cite the source of borrowed images.
\item Check readability from the back of the room: no chart should need a magnifying glass.
\end{enumerate}
\textbf{Reflex:} a chart replaces a paragraph of numbers; choose the chart type \emph{from the nature of the data}, not from taste.
\end{methode}

\begin{aretenir}[Choosing the right chart type]
\begin{center}
\begin{tabularx}{\linewidth}{|l|X|X|}
\hline
\rowcolor{popblueL} \textbf{Chart type} & \textbf{Nature of the data} & \textbf{Example} \\
\hline
Column / bar & Comparing categories at one moment & Marks of 5 subjects \\
\hline
Line & Evolution of a quantity over time & Subscribers 2018--2023 \\
\hline
Pie / doughnut & Shares (percentages) of a total & Sources of internet access \\
\hline
Table & Exact numerical values & Price list in FCFA \\
\hline
SmartArt & Process, hierarchy, cycle & Water cycle, staff organigram \\
\hline
\end{tabularx}
\end{center}
\end{aretenir}

\begin{exoresolu}[Choosing the right illustration]
\textbf{Statement.} A group presents a survey of 200 students of a school in Yaound\'e on their main source of internet access: mobile data 130, home Wi-Fi 40, cybercaf\'e 20, no access 10. (a) Which chart type is most appropriate to show the \emph{shares} of each source? Justify. (b) Which chart type would be appropriate to show how the number of internet users in the school evolved from 2020 to 2024? (c) The group wants to show the exact figures in a structured way on one slide: which object should they use? (d) State one precaution when inserting a photo downloaded from the internet.
\tcblower
\textbf{Solution.}
\begin{enumerate}
\item[(a)] A \cle{pie chart} (or doughnut chart): the four categories are parts of a whole (200 students), and a pie chart displays percentages of a total. The slices would represent $65\ \%$, $20\ \%$, $10\ \%$ and $5\ \%$ respectively (e.g. $130/200 = 0.65 = 65\ \%$).
\item[(b)] A \cle{line chart}: it is the standard choice to show the \emph{evolution of a quantity over time} (years 2020--2024 on the horizontal axis, number of users on the vertical axis).
\item[(c)] A \cle{table}: it presents exact numerical values in rows and columns (Source | Number of students | Percentage), which a chart cannot show precisely.
\item[(d)] Precautions (any one): check that the image is free of rights or cite the \emph{source}; verify the image is of good \emph{resolution} so it does not appear blurred when projected; compress it so the file stays light.
\end{enumerate}
\textbf{Examiner's remark:} the justification must refer to the \emph{nature of the data} (parts of a whole $\rightarrow$ pie; time series $\rightarrow$ line), which is what earns full marks.
\end{exoresolu}

\coursec{Multimedia, Transitions and Animations}

\courssub{Skill 4: Applying transitions and animations}

\begin{definition}[Transition vs animation]
A \cle{transition} is a visual effect applied \emph{between two slides}, i.e. the way slide $n+1$ replaces slide $n$ on the screen. An \cle{animation} is an effect applied to an \emph{object inside a slide} (text, picture, chart): the four families are \emph{Entrance} (the object appears), \emph{Emphasis} (the object is highlighted), \emph{Exit} (the object disappears) and \emph{Motion paths} (the object moves).
\end{definition}

\begin{methode}[Using transitions and animations correctly]
\begin{enumerate}
\item Distinguish the two concepts: a \cle{transition} acts \emph{between slides}; an \cle{animation} acts \emph{on objects within a slide}.
\item \textbf{Transitions:} select the slide(s), \emph{Transitions} tab, choose a sober effect (Fade, Push), set the \emph{duration}, and choose the advance mode: \emph{On mouse click} or \emph{After} a given time. Use \emph{Apply to All} for uniformity.
\item \textbf{Animations:} select the object, \emph{Animations} tab, then pick an Entrance, Emphasis, Exit or Motion path effect.
\item Order and time the animations in the \emph{Animation Pane}: start \emph{On Click}, \emph{With Previous} or \emph{After Previous}.
\item \textbf{Sobriety rule:} one transition for the whole show, at most 1--2 animation effects per slide, each serving a pedagogical purpose (e.g. revealing bullet points one by one).
\end{enumerate}
\end{methode}

\begin{attention}[The classic exam trap]
Do not confuse \emph{transition} and \emph{animation}. ``The slide appears with a fade'' is a \textbf{transition}; ``the bullet points appear one by one'' is an \textbf{animation}. Mixing them up in a GCE answer loses the mark even if the rest is correct.
\end{attention}

\begin{exoresolu}[Planning transitions and animations]
\textbf{Statement.} For a presentation on \emph{The fight against malaria in Cameroon}, a student wants: (i) all slides to appear with the same discreet effect lasting 1 second; (ii) on slide 4, the three symptoms of malaria to appear one after another, each when the presenter clicks; (iii) on slide 5, a mosquito picture to move across the screen automatically after the title appears. (a) For each wish, state whether it is a transition or an animation, and give the exact settings. (b) The student applies a different, spectacular transition to each of his 15 slides. Criticise this choice.
\tcblower
\textbf{Solution.}
\begin{enumerate}
\item[(a)] (i) \textbf{Transition:} select all slides (\texttt{Ctrl+A} in the slide pane), \emph{Transitions} tab, choose \emph{Fade}, set \emph{Duration} $= 1.00$ s, then click \emph{Apply to All}. (ii) \textbf{Animation (Entrance):} select the bulleted list of symptoms on slide 4, apply an Entrance effect (e.g. \emph{Appear} or \emph{Fade}); in the Animation Pane, set each bullet to start \emph{On Click}, so the three symptoms appear successively at each click. (iii) \textbf{Animations combined:} give the title an Entrance effect, then select the mosquito picture and apply a \emph{Motion path} (e.g. a line across the slide); in the Animation Pane set the motion path to start \emph{After Previous} so it plays automatically once the title has appeared.
\item[(b)] Using 15 different spectacular transitions is a mistake: it \emph{distracts the audience} from the content, looks unprofessional, slows the presentation and can tire the viewers. Good practice is a single sober transition (or none) applied to all slides, with effects reserved for moments where they support understanding.
\end{enumerate}
\textbf{Examiner's remark:} clearly separating \emph{transition = between slides} and \emph{animation = on objects} is the core marking point; confusing them loses marks.
\end{exoresolu}

\courssub{Skill 5: Inserting sound and video (multimedia)}

\begin{methode}[Making a presentation multimedia]
\begin{enumerate}
\item \textbf{Insert audio:} \emph{Insert $\rightarrow$ Audio} (from file). Choose the start mode: \emph{Automatically}, \emph{On Click}, or \emph{Play across slides} for background music.
\item \textbf{Insert video:} \emph{Insert $\rightarrow$ Video}; prefer short clips (a few dozen seconds) in common formats (\texttt{.mp4}, \texttt{.wmv}); trim the clip with the \emph{Trim Video} tool.
\item \textbf{Keep media with the file:} copy the media files into the same folder as the presentation \emph{before} inserting, or embed them; otherwise the sound/video will not play on another computer.
\item \textbf{Test on the presentation computer} (the school projector room), not only on your own machine, and check the loudspeakers.
\item Respect \emph{copyright}: use your own recordings or royalty-free media, and cite sources.
\item Balance: multimedia must \emph{illustrate} the message (e.g. a 30 s video of a mobile money transaction in a talk on e-commerce), never replace it.
\end{enumerate}
\end{methode}

\begin{savaistu}[Linked or embedded media]
When you insert a sound or video, the software may only store a \emph{link} (the path to the file on your disk) instead of \emph{embedding} a copy inside the presentation. This is why a presentation that works perfectly at home can fail on the school computer: the linked file was left behind. Always transport the media files together with the \texttt{.pptx} file, or embed them, and test on the target machine.
\end{savaistu}

\begin{exoresolu}[Embedding multimedia correctly]
\textbf{Statement.} Marie prepares a presentation on \emph{Traditional dances of the West Region}. She inserts a 3-minute video filmed at a festival in Bafoussam and a background song. (a) At school, on the classroom computer, the video refuses to play and an error message appears. Give two likely causes and the remedy for each. (b) She wants the song to continue while she moves through slides 2 to 6: which option must she activate? (c) The 3-minute video interrupts her talk for too long. Propose two solutions. (d) State one legal precaution concerning the song.
\tcblower
\textbf{Solution.}
\begin{enumerate}
\item[(a)] Likely causes and remedies: (i) the video file was only \emph{linked}, not embedded, and was not copied to the school computer $\Rightarrow$ copy the video file together with the presentation (same folder) or embed it, then re-insert; (ii) the classroom computer lacks the \emph{codec}/format support $\Rightarrow$ convert the video to a widely supported format such as \texttt{.mp4} before the presentation. (A third acceptable cause: the file path changed because the presentation was moved.)
\item[(b)] She must set the audio to \emph{Play across slides} (in the Playback options of the audio object), with start \emph{Automatically} on slide 2 and stop after slide 6 if the tool allows it.
\item[(c)] Solutions: (i) use \emph{Trim Video} to keep only the most illustrative 30--40 seconds; (ii) split the video and show short extracts at the relevant moments, commenting over them, instead of one long uninterrupted clip.
\item[(d)] The song is protected by \emph{copyright}: she should use a royalty-free track, obtain authorisation, or at least cite the artist and source, especially if the presentation will be published online.
\end{enumerate}
\textbf{Examiner's remark:} ``linked vs embedded media'' and ``test on the target machine'' are the classic expected answers for multimedia problems.
\end{exoresolu}

\coursec{Delivering, Printing and Evaluating the Presentation}

\courssub{Skill 6: Delivering, printing and exporting the presentation}

\begin{methode}[Presenting, printing handouts and exporting]
\begin{enumerate}
\item \textbf{Rehearse} with \emph{Rehearse Timings} and check the total duration against the allowed time.
\item Run the show with \texttt{F5} (from the beginning) or \texttt{Shift+F5} (from the current slide); navigate with arrow keys, \texttt{B} to black the screen, \texttt{Esc} to quit.
\item Use \emph{Presenter View} (speaker notes visible to you, slides only to the audience) when available; write \emph{speaker notes} under each slide.
\item \textbf{Print handouts:} \emph{File $\rightarrow$ Print $\rightarrow$ Handouts} (2, 3, 4, 6 or 9 slides per page); the 3-slide layout leaves lines for the audience's notes. Print \emph{Notes Pages} for yourself.
\item \textbf{Export} a copy as \cle{PDF} to guarantee identical display on any computer, and keep the editable \texttt{.pptx} file.
\item Before the day: copy everything (presentation + media + PDF) on a USB key and test on the actual projector (resolution, sound, cables).
\end{enumerate}
\end{methode}

\begin{exoresolu}[Preparing the delivery]
\textbf{Statement.} A candidate must defend a project in 10 minutes before a jury, using a 12-slide presentation. (a) Explain how he can check, before the jury day, that his speech fits in 10 minutes. (b) The jury asks for a paper document containing all the slides, four per page, to annotate during the defence. Describe how to produce it. (c) On the jury's computer, the fonts of his presentation are replaced and the layout is broken. What should he have done to avoid this? (d) Give two keyboard shortcuts useful during the slide show and their effects.
\tcblower
\textbf{Solution.}
\begin{enumerate}
\item[(a)] He uses \emph{Slide Show $\rightarrow$ Rehearse Timings}: the software records the time spent on each slide while he practises aloud, and displays the total. If the total exceeds 10 minutes, he shortens his comments or removes a slide, then rehearses again until the timing is correct. He can also save these timings for an automatic show.
\item[(b)] \emph{File $\rightarrow$ Print}; under Settings, choose \emph{Handouts (4 slides per page)} (horizontal or vertical order), select the printer and print. The jury then receives every slide in miniature with space around for annotations. (The 3-slides-per-page layout, with ruled lines, is also acceptable if justified.)
\item[(c)] He should have \emph{exported a PDF copy} of the presentation (\emph{File $\rightarrow$ Export/Save As $\rightarrow$ PDF}): a PDF preserves fonts, images and layout exactly on any computer. Alternatively, embedding fonts when saving, or using only standard fonts, reduces the risk.
\item[(d)] Any two: \texttt{F5} starts the show from the first slide; \texttt{Shift+F5} starts from the current slide; \texttt{B} (or \texttt{W}) temporarily blacks (whites) the screen to refocus attention on the speaker; \texttt{Esc} ends the show; arrow keys/\texttt{Space} advance to the next slide.
\end{enumerate}
\textbf{Examiner's remark:} in (c), ``export to PDF'' is the strongest answer because it addresses \emph{both} fonts and layout.
\end{exoresolu}

\courssub{Skill 7: Evaluating and improving a presentation (integration)}

\begin{methode}[Evaluating a slide presentation with a grid]
To evaluate (or self-evaluate) a presentation, use a criteria grid covering four families:
\begin{enumerate}
\item \textbf{Content:} accurate information, clear structure (introduction, development, conclusion), one idea per slide, adapted to the audience and to the time limit.
\item \textbf{Design:} consistent theme/master, readable fonts and sizes, high contrast, limited colours, well-chosen illustrations (correct chart types), no overcrowding.
\item \textbf{Multimedia and effects:} sober transitions, purposeful animations, working sound/video, embedded media.
\item \textbf{Delivery and technical quality:} correct spelling, speaker notes, rehearsed timing, tested on the target machine, PDF/USB backup.
\end{enumerate}
For each criterion, tick \emph{Yes/Partially/No}, then list concrete improvements. This is exactly the approach expected in the \emph{integration activities} lesson: produce, present, then critique and improve.
\end{methode}

\begin{exoresolu}[Integration: critiquing and correcting a presentation]
\textbf{Statement.} Here is the description of a classmate's 8-slide presentation on \emph{Mobile Money in Cameroon}: slide 1 has no title; the text is written in long paragraphs in 14 pt yellow font on a white background; each slide uses a different theme; a pie chart is used to show the evolution of the number of MTN Mobile Money subscribers from 2018 to 2023; every slide has a different spinning transition; a video is inserted but does not play on the school computer; the file is named \texttt{doc1.pptx}. Identify \textbf{six} weaknesses and, for each, give the correction.
\tcblower
\textbf{Solution.} Any six of the following (weakness $\rightarrow$ correction):
\begin{enumerate}
\item \textbf{Missing title on slide 1} $\rightarrow$ use the \emph{Title Slide} layout with the presentation title, the author's name, class and date.
\item \textbf{Long paragraphs} $\rightarrow$ replace them by short bulleted keywords (about 6 lines per slide); the details belong to the speech and the speaker notes.
\item \textbf{Unreadable text (14 pt, yellow on white)} $\rightarrow$ increase body text to at least 24 pt and use a \emph{high-contrast} combination, e.g. dark text on a light background.
\item \textbf{A different theme per slide} $\rightarrow$ apply a \emph{single theme} (or the Slide Master) to the whole presentation for visual consistency.
\item \textbf{Wrong chart type} $\rightarrow$ an evolution over time (2018--2023) must be shown with a \emph{line chart} (or column chart); a pie chart is only for shares of a total at one moment.
\item \textbf{Multiple spinning transitions} $\rightarrow$ apply \emph{one sober transition} (e.g. Fade) to all slides; effects must serve the message, not distract.
\item \textbf{Video not playing} $\rightarrow$ copy the video file with the presentation (or embed it), use a common format (\texttt{.mp4}) and \emph{test on the school computer} beforehand.
\item \textbf{Meaningless file name} $\rightarrow$ save as e.g. \texttt{MobileMoney\_Presentation.pptx} and keep a PDF export plus a USB backup.
\end{enumerate}
\textbf{Conclusion:} after correction, the presentation satisfies the four evaluation families: clear content, consistent and readable design, sober effects, and technical reliability.
\textbf{Examiner's remark:} in integration questions, always pair \emph{each weakness with its concrete remedy}; a list of faults without corrections earns only half the marks.
\end{exoresolu}

\begin{exercice}[Practice: from production to delivery]
Your class must present \emph{``The role of ANTIC in securing online services in Cameroon''} in 8 minutes.
\begin{enumerate}
\item Propose a plan of 6 slides, giving for each slide its \emph{title} and its \emph{layout}.
\item State the chart type you would use to show the growth of reported cyber-incidents over five years, and justify.
\item You want a short recorded interview (audio) to play only when you click on slide 4. Give the exact setting.
\item Give three checks to perform on the school computer the day before the presentation.
\end{enumerate}
\end{exercice}

\begin{corrige}[Practice: from production to delivery]
\begin{enumerate}
\item Example plan: slide 1 \emph{Title Slide} (title, names, class); slide 2 \emph{Title and Content} (What is ANTIC? missions); slide 3 \emph{Title and Content} (services: certificates, awareness); slide 4 \emph{Title and Content} (recorded interview + audio icon); slide 5 \emph{Title and Content} (chart of cyber-incidents); slide 6 \emph{Title and Content} or \emph{Blank} (conclusion and thanks).
\item A \emph{line chart} (or column chart): the data describe an \emph{evolution over time}, which is exactly what a line chart shows; a pie chart would be wrong because the years are not parts of a whole.
\item Select the audio object on slide 4, open the \emph{Playback} options and set \emph{Start: On Click} (not Automatically, not Play across slides).
\item Any three: check that the presentation file opens; check that the audio/video media play (embedded or copied with the file); check the projector display and resolution; check the loudspeakers; check that fonts and layout are correct (or use the PDF copy); check that the USB key is readable.
\end{enumerate}
\end{corrige}

\begin{aretenir}[Chapter summary]
\begin{itemize}
\item Plan first: \emph{one slide = one message}; choose the \emph{layout} before typing; save often under a meaningful name.
\item Consistency comes from the \emph{theme} and the \emph{Slide Master}; manual formatting overrides the master (use \emph{Reset} to fix).
\item Choose illustrations from the \emph{nature of the data}: pie = shares of a total, line = evolution over time, column = comparison, table = exact values.
\item \emph{Transition} = between slides; \emph{animation} = on objects (Entrance, Emphasis, Exit, Motion path). Stay sober: one transition, purposeful animations.
\item Multimedia: embed or transport media files, use common formats (\texttt{.mp4}), respect copyright, and \emph{test on the target machine}.
\item Deliver: rehearse timings, use speaker notes and Presenter View, print handouts, export a \emph{PDF} backup, and evaluate with a four-family criteria grid.
\end{itemize}
\end{aretenir}

\newpage

\coursec{Exercises}

\courssub{I check my knowledge}

\begin{exercice}[MCQ: transitions or animations?]
For each of the following effects, state whether it is a \textbf{transition} or an \textbf{animation}. Choose the correct letter.
\begin{enumerate}
\item A ``Fade'' effect that appears when moving from Slide 3 to Slide 4.
\item A ``Fly In'' effect applied to the title text of Slide 2.
\item A ``Wipe'' effect that makes a chart appear bar by bar on Slide 5.
\item A ``Morph'' effect that transforms the whole of Slide 6 into Slide 7.
\end{enumerate}
\textbf{Choices:}
\begin{enumerate}
\item[a)] 1 = transition, 2 = animation, 3 = animation, 4 = transition
\item[b)] 1 = animation, 2 = transition, 3 = transition, 4 = animation
\item[c)] All four are transitions.
\item[d)] All four are animations.
\end{enumerate}
\end{exercice}

\begin{exercice}[True or false — justify]
State whether each statement is \textbf{true} or \textbf{false}, and justify your answer in one or two sentences.
\begin{enumerate}
\item The Slide Master allows you to place the school logo and slide numbers on every slide at once.
\item Transitions are effects applied to individual objects (text, images, charts) on a slide.
\item A presentation saved with a video that was only \emph{linked} (not embedded) will always play correctly on any other computer.
\item The 6$\times$6 rule recommends at most 6 bullet points per slide and at most 6 words per bullet point.
\end{enumerate}
\end{exercice}

\begin{exercice}[Key definitions]
Define each of the following terms in one or two clear sentences, as expected at the GCE Advanced Level:
\begin{enumerate}
\item Slide Master
\item Transition
\item Animation
\item Handout
\item Slide Show (presentation mode)
\end{enumerate}
\end{exercice}

\begin{exercice}[Direct calculation: timing a slide show]
A student in Upper Sixth prepares a presentation of 12 slides. The slide show is set to advance \textbf{automatically}. Slides 1 to 4 advance after 8 seconds each, slides 5 to 10 after 15 seconds each, and slides 11 and 12 after 20 seconds each.
\begin{enumerate}
\item Calculate the total duration of the slide show in seconds.
\item Convert this duration into minutes and seconds.
\item The student must present within a 3-minute limit. By how many seconds must the total timing be reduced?
\end{enumerate}
\end{exercice}

\courssub{I practise}

\begin{exercice}[Restructuring a raw presentation]
Your classmate gives you a raw 8-slide presentation about ``The History of Our School''. The fonts change from slide to slide, there is no logo, no slide numbers, and the slides are in a random order (the conclusion appears as Slide 3). Describe, step by step, the operations you would perform in the presentation software to:
\begin{enumerate}
\item apply a single consistent theme to all slides;
\item use the Slide Master to display the school logo and automatic slide numbers on every slide;
\item reorder the slides into a logical sequence (title, introduction, body, conclusion).
\end{enumerate}
\end{exercice}

\begin{exercice}[Inserting and animating a chart]
A school wants to show its GCE Advanced Level pass rates: 72\% in 2022, 78\% in 2023 and 85\% in 2024.
\begin{enumerate}
\item Describe how to insert a \textbf{column chart} of these three values on a slide.
\item Explain how to animate the chart \emph{series by series} (or category by category) so that each year's column appears separately.
\item Explain how to set this slide to advance automatically after 10 seconds.
\end{enumerate}
\end{exercice}

\begin{exercice}[Producing handouts and notes pages]
You must give a printed document to an audience of 30 parents during a school open day in Bamenda.
\begin{enumerate}
\item Describe the steps to print \textbf{handouts with 6 slides per page}, leaving lines or space where the audience can write notes.
\item Explain in which situations \textbf{notes pages} (one slide plus the speaker's notes) are preferred to handouts.
\end{enumerate}
\end{exercice}

\begin{exercice}[Evaluating two slides]
Two slides are presented to you. \textbf{Slide A} contains 12 bullet points of long sentences, yellow text on a white background, three different fonts, a low-quality stretched image and five different animations. \textbf{Slide B} contains 5 short bullet points, dark blue text on a light background, one font, one clear image and one simple transition.
\begin{enumerate}
\item Identify \textbf{five design faults} in Slide A.
\item For each fault, propose a correction, referring to the 6$\times$6 rule and to the principles of contrast and consistency.
\end{enumerate}
\end{exercice}

\courssub{I prepare for the GCE}

\begin{exercice}[Structured question: planning a business pitch (8 marks)]
A young entrepreneur in Douala must present a business plan to a bank in order to obtain a loan. He decides to use presentation software.
\begin{enumerate}
\item List and explain the \textbf{five steps} for planning an effective presentation, from defining the objective to rehearsing the delivery. \hfill (5 marks)
\item Give \textbf{three reasons} why the Slide Master is useful in preparing this presentation. \hfill (3 marks)
\end{enumerate}
\end{exercice}

\begin{exercice}[Scenario question: the video that does not play (6 marks)]
During a GCE practice presentation, a student prepares her slides at home and inserts a video stored in a folder on her personal computer. At school, the video refuses to play on the classroom computer, although all the other slides display correctly.
\begin{enumerate}
\item Explain the most likely cause of this problem, using the terms \emph{linked file} and \emph{embedded file}. \hfill (2 marks)
\item Describe \textbf{two different methods} the student could have used to prevent this problem. \hfill (4 marks)
\end{enumerate}
\end{exercice}

\begin{exercice}[GCE-style essay (10 marks)]
\textbf{``Animations and transitions always improve a presentation.''}

Discuss this statement. Your essay should:
\begin{enumerate}
\item explain what transitions and animations are, with one example of each;
\item present at least \textbf{two arguments in favour} of using them (with examples, e.g.\ revealing bullet points progressively, showing a chart series by series);
\item present at least \textbf{two arguments against} their excessive use (with examples, e.g.\ distraction, loss of time, unprofessional appearance);
\item conclude with clear \textbf{guidelines} for the reasonable use of transitions and animations in an academic or professional presentation.
\end{enumerate}
Your answer should be well structured (introduction, development, conclusion) and illustrated with concrete examples. \hfill (10 marks)
\end{exercice}

\newpage

\coursec{Solutions to the exercises}

\courssub{Solutions — I check my knowledge}

\begin{corrige}[Exercise 1 — MCQ: transitions or animations?]
The correct answer is \textbf{a)}: 1 = transition, 2 = animation, 3 = animation, 4 = transition.

\textbf{Justification.} A \cle{transition} is a visual effect that occurs \emph{between two slides}, i.e.\ when the slide show moves from one slide to the next. An \cle{animation} is a visual effect applied to an \emph{object inside a single slide} (text, picture, chart, shape).
\begin{enumerate}
\item The ``Fade'' appears \textbf{when moving from Slide 3 to Slide 4}: it concerns the passage between two slides $\Rightarrow$ \textbf{transition}.
\item The ``Fly In'' is applied to the \textbf{title text}, an object on Slide 2 $\Rightarrow$ \textbf{animation}.
\item The ``Wipe'' makes the \textbf{chart} (an object on Slide 5) appear bar by bar $\Rightarrow$ \textbf{animation}.
\item ``Morph'' transforms \textbf{the whole of Slide 6 into Slide 7}: it acts between two slides $\Rightarrow$ \textbf{transition}.
\end{enumerate}
\begin{attention}[The classic trap]
Do not be misled by the name of the effect (``Fade'', ``Wipe'' exist in both families). The only reliable criterion is the \textbf{target}: whole slide / passage between slides = transition; object inside a slide = animation.
\end{attention}
\end{corrige}

\begin{corrige}[Exercise 2 — True or false]
\begin{enumerate}
\item \textbf{True.} The Slide Master controls the layout of all slides: any element placed on it (logo, slide number field, footer, background) is automatically reproduced on every slide using that master, in a single operation.
\item \textbf{False.} Effects applied to individual objects are \emph{animations}. Transitions apply to the passage \emph{between two slides}, not to objects.
\item \textbf{False.} A linked video is not copied into the presentation file; the file only stores the \emph{path} to the video. If the video is not copied with the presentation (or the path changes), it will not play on another computer. Only an \emph{embedded} video travels with the file.
\item \textbf{True.} The 6$\times$6 rule is a readability guideline: maximum 6 bullet points per slide and maximum 6 words per bullet point, so that slides remain airy and the audience listens to the speaker instead of reading.
\end{enumerate}
\end{corrige}

\begin{corrige}[Exercise 3 — Key definitions]
\begin{enumerate}
\item \textbf{Slide Master:} the model slide that stores the common design (theme, fonts, colours, logo, footer, slide numbers) of a presentation. Any change made on the master is automatically applied to all the slides that depend on it.
\item \textbf{Transition:} a visual effect that plays when the slide show moves from one slide to the next (e.g.\ Fade, Morph, Push).
\item \textbf{Animation:} a visual effect applied to an object (text, image, chart) within a slide, controlling how and when it appears, is emphasised or disappears (e.g.\ Fly In, Wipe).
\item \textbf{Handout:} a printed version of the presentation showing several miniatures of slides per page (typically 2, 3, 4, 6 or 9), given to the audience so they can follow and take notes.
\item \textbf{Slide Show (presentation mode):} the full-screen playback mode in which the presentation is displayed to the audience, with its transitions, animations and timings, starting from the current or the first slide.
\end{enumerate}
\end{corrige}

\begin{corrige}[Exercise 4 — Direct calculation: timing a slide show]
\begin{methode}[Calculating total slide show duration]
\begin{itemize}
\item List each group of slides with its individual duration.
\item Multiply the number of slides by the duration per slide for each group.
\item Sum all the products to obtain the total duration in seconds.
\item Convert seconds to minutes and seconds (divide by 60).
\item Compare with the time limit to decide if reduction is needed.
\end{itemize}
\end{methode}
\begin{enumerate}
\item Total duration:
\begin{align*}
T &= \underbrace{4 \times 8}_{\text{slides 1--4}} + \underbrace{6 \times 15}_{\text{slides 5--10}} + \underbrace{2 \times 20}_{\text{slides 11--12}}\\
&= 32 + 90 + 40 = 162\ \text{s}.
\end{align*}
\item Conversion: $162 = 2 \times 60 + 42$, so the show lasts \textbf{2 minutes 42 seconds}.
\item The limit is 3 minutes $= 180$ s. Since $162 < 180$, the show is already \emph{within} the limit: it must be reduced by $162 - 180 = -18$ s, i.e.\ \textbf{no reduction is needed}; the student even has \textbf{18 seconds of margin}.
\end{enumerate}
\begin{attention}[Read before computing]
Question 3 is a trap: many candidates compute $180 - 162 = 18$ and answer ``reduce by 18 s''. Always compare first: here the timing is \emph{below} the limit, so nothing must be removed.
\end{attention}
\end{corrige}

\courssub{Solutions — I practise}

\begin{corrige}[Exercise 5 — Restructuring a raw presentation]
\begin{enumerate}
\item \textbf{Apply one consistent theme.}
Open the \textbf{Design} tab, choose a single theme in the Themes gallery and click it: it is applied to the whole presentation at once (fonts, colours, background). If needed, use the \emph{Variants} group to adjust the colour scheme and font set. This immediately removes the mixture of fonts.
\item \textbf{Add logo and slide numbers via the Slide Master.}
Go to \textbf{View $\rightarrow$ Slide Master}. Select the top (parent) master so the change affects all layouts. Use \textbf{Insert $\rightarrow$ Pictures} to place the school logo (e.g.\ top-right corner), then \textbf{Insert $\rightarrow$ Slide Number} (or Header \& Footer, tick \emph{Slide number}, then \emph{Apply to All}). Close the Master view: the logo and automatic numbers now appear on every slide.
\item \textbf{Reorder the slides.}
Switch to \textbf{View $\rightarrow$ Slide Sorter} (or the left thumbnail pane in Normal view). Drag and drop the thumbnails into a logical order: (1) title slide, (2) introduction/outline, (3--6) body in chronological order of the school's history, (7) conclusion, (8) references/thank-you. Move the conclusion from position 3 to the end. Finally, run the slide show from the beginning to check the new sequence.
\end{enumerate}
\begin{savaistu}[Cameroon context]
In many Cameroonian schools (e.g.\ GBHS Bamenda, Lycée Leclerc Yaoundé), the ICT practical exam requires candidates to demonstrate Slide Master usage. A common task: insert the school logo (often the coat of arms) and automatic slide numbers on all slides in under two minutes — only the Master makes this possible.
\end{savaistu}
\end{corrige}

\begin{corrige}[Exercise 6 — Inserting and animating a chart]
\begin{enumerate}
\item \textbf{Inserting the column chart.} On the slide, go to \textbf{Insert $\rightarrow$ Chart}, choose \textbf{Column $\rightarrow$ Clustered Column} and validate. A small data sheet opens: enter the categories 2022, 2023, 2024 and the values 72, 78, 85 (as one series ``Pass rate (\%)''). Close the data sheet; add a chart title such as ``GCE A-Level pass rate''.
\item \textbf{Animating series by series / category by category.} Select the chart, open the \textbf{Animations} tab and choose an entrance effect (e.g.\ \emph{Wipe} or \emph{Appear}). Then open \textbf{Effect Options} and set \emph{Sequence} to \textbf{By Category} (or ``By Series''): each year's column now appears separately, on a click or automatically after the previous one. This lets the speaker comment on each year as it appears.
\item \textbf{Automatic advance after 10 s.} Select the slide, go to the \textbf{Transitions} tab, in the \emph{Timing} group untick \emph{On Mouse Click} (optional), tick \textbf{After} and enter \textbf{00:10.00}. The slide will advance automatically 10 seconds after it appears.
\end{enumerate}
\begin{attention}[Chart animation trap]
If you choose ``As One Object'' the whole chart appears at once — no progressive reveal. Always select \emph{By Category} or \emph{By Series} in Effect Options to animate data elements individually.
\end{attention}
\end{corrige}

\begin{corrige}[Exercise 7 — Producing handouts and notes pages]
\begin{enumerate}
\item \textbf{Handouts, 6 slides per page.} Go to \textbf{File $\rightarrow$ Print}. In the Settings section, open the layout menu (``Full Page Slides'' by default) and under \emph{Handouts} choose \textbf{6 Slides Horizontal} (or Vertical). This layout leaves free space around the miniatures where parents can write; choosing the \textbf{3 Slides} layout instead gives ruled lines next to each slide. Set the number of copies to 30, check the preview, then print.
\item \textbf{When notes pages are preferred.} Notes pages (one slide at the top, the speaker's notes below) are intended for the \emph{presenter}, not the audience. They are preferred when: rehearsing the talk; the speech must be delivered by someone else (a colleague replaces the speaker); the presentation is archived as a complete document; or detailed explanations, figures and sources must accompany each slide without overloading the slide itself. Handouts remain the right choice for the audience.
\end{enumerate}
\begin{savaistu}[Printing in Cameroon]
Many schools print handouts on shared network printers (often managed by CAMTEL or MTN Business solutions). Always preview in \textbf{Print Preview} to avoid wasting paper — a real concern when printing 30 copies for a PTA meeting at a rural school.
\end{savaistu}
\end{corrige}

\begin{corrige}[Exercise 8 — Evaluating two slides]
\begin{enumerate}
\item \textbf{Five faults of Slide A:} (i) 12 long bullet points — information overload; (ii) yellow text on white background — insufficient contrast, unreadable; (iii) three different fonts — inconsistency; (iv) a stretched, low-quality image — distorted and unprofessional; (v) five different animations — distracting visual noise.
\item \textbf{Corrections:}
\begin{center}\begin{tabularx}{\linewidth}{|l|X|}\hline
\rowcolor{popblueL} \textbf{Fault} & \textbf{Correction} \\ \hline
12 long bullets & Apply the 6$\times$6 rule: at most 6 bullets of at most 6 words; split the content over two slides and let the speech carry the details. \\ \hline
Yellow on white & Respect contrast: dark text on a light background (or the reverse), e.g.\ dark blue on white as in Slide B. \\ \hline
Three fonts & Consistency: one single font (at most two: one for titles, one for body), fixed in the Slide Master. \\ \hline
Stretched image & Insert a good-resolution image and resize it by the corner handles only, keeping the proportions (aspect ratio). \\ \hline
Five animations & Keep at most one simple, purposeful animation (e.g.\ reveal bullets progressively) and one discreet transition for the whole deck. \\ \hline
\end{tabularx}\end{center}
Slide B follows all these principles: it is the model to imitate.
\end{enumerate}
\end{corrige}

\courssub{Solutions — I prepare for the GCE}

\begin{corrige}[Exercise 9 — Structured question: planning a business pitch]
\textbf{Indicative mark scheme (8 marks).} Question 1: 1 mark per correctly named \emph{and} explained step ($5 \times 1$). Question 2: 1 mark per valid reason ($3 \times 1$). Naming without explanation earns at most half the marks.

\textbf{1. The five steps of planning (5 marks):}
\begin{enumerate}
\item \textbf{Define the objective:} decide what the presentation must achieve — here, convince the bank to grant the loan. Every slide must serve this goal.
\item \textbf{Analyse the audience:} bankers expect figures, reliability and seriousness; the content, vocabulary and sober design are adapted accordingly.
\item \textbf{Plan and structure the content:} build an outline — introduction, presentation of the business, market study, financial plan, loan requested, conclusion — and select only the essential points.
\item \textbf{Design and produce the slides:} apply a consistent theme and Slide Master, respect the 6$\times$6 rule, insert clear charts (e.g.\ projected revenue) rather than long text.
\item \textbf{Rehearse and check the delivery:} run the slide show, time it, check transitions, media and spelling, and practise speaking with the slides until the timing fits the appointment.
\end{enumerate}

\textbf{2. Three reasons the Slide Master is useful (3 marks):}
\begin{enumerate}
\item \textbf{Consistency:} it guarantees the same fonts, colours and layout on all slides, giving a professional image to the bank.
\item \textbf{Time saving:} the logo, footer and slide numbers are inserted \emph{once} on the master instead of slide by slide.
\item \textbf{Easy global modification:} changing the design (e.g.\ the colour scheme) on the master updates the whole presentation instantly, avoiding inconsistencies.
\end{enumerate}
\textbf{Examiner's expectations:} structured answer, correct vocabulary (objective, audience, outline, master), and reasons clearly linked to the banking context.
\end{corrige}

\begin{corrige}[Exercise 10 — Scenario question: the video that does not play]
\textbf{Indicative mark scheme (6 marks).} (a) 1 mark for identifying linking as the cause, 1 mark for the explanation with correct terms. (b) 2 marks per method (1 for the method, 1 for the description), $2 \times 2 = 4$.

\textbf{1. Most likely cause (2 marks).} At home the student inserted the video as a \emph{linked file}: the presentation only stores the \emph{path} to the video on her personal computer (e.g.\ \texttt{C:\textbackslash Users\textbackslash...\textbackslash video.mp4}), not the video itself. On the classroom computer this path does not exist, so the link is broken and the video cannot play, while the rest of the slides — which are embedded in the file — display normally. An \emph{embedded file}, by contrast, is physically copied inside the presentation and travels with it.

\textbf{2. Two prevention methods (4 marks).}
\begin{enumerate}
\item \textbf{Embed the video:} insert it so that it is copied into the presentation file (or use \emph{Insert $\rightarrow$ Video} then save in a format that embeds media). The file becomes heavier, but it is self-contained and plays on any computer.
\item \textbf{Copy the video together with the presentation, keeping the link valid:} place the video in the \emph{same folder} as the presentation \emph{before} inserting it, then copy the whole folder (e.g.\ on a USB flash drive) to the school computer without changing the relative paths. One can also use the \emph{Package/Export for CD} feature, which gathers the presentation and all linked media automatically.
\end{enumerate}
\textbf{Examiner's expectations:} the pair \emph{linked/embedded} must be used correctly; any two coherent, technically valid methods are accepted.
\begin{savaistu}[ANTIC recommendation]
The National Agency for Information and Communication Technologies (ANTIC) recommends embedding media for official presentations (e.g.\ MINEDUB reports) to avoid broken links during inter-ministerial meetings where USB drives move between different Windows versions.
\end{savaistu}
\end{corrige}

\begin{corrige}[Exercise 11 — GCE-style essay]
\textbf{Indicative mark scheme (10 marks):} definitions with examples 2 marks; two arguments in favour with examples 2 marks; two arguments against with examples 2 marks; guidelines 2 marks; structure (introduction, development, conclusion) and quality of expression 2 marks.

\textbf{Model answer.}

\emph{Introduction.} Modern presentation software offers two families of effects: \textbf{transitions}, which animate the passage from one slide to the next (e.g.\ a \emph{Fade} between two slides), and \textbf{animations}, which animate objects inside a slide (e.g.\ a title that \emph{flies in}). The claim that they ``always improve'' a presentation deserves discussion: used with measure they help, used in excess they harm.

\emph{Arguments in favour.} First, animations \textbf{control the flow of information}: revealing bullet points one by one keeps the audience focused on the point being discussed instead of reading ahead. Secondly, they can \textbf{clarify data}: animating a chart series by series (for example the GCE pass rates year by year) lets the speaker comment on each value as it appears. A discreet transition, finally, gives rhythm and signals a change of section.

\emph{Arguments against excess.} First, too many effects \textbf{distract}: five different animations on one slide draw attention to the effects rather than to the message. Secondly, they \textbf{waste time} and can look \textbf{unprofessional}: a bank manager or an examiner faced with spinning text and sound effects will doubt the seriousness of the speaker; long effects also lengthen the show and may break the timing.

\emph{Guidelines.} Use at most \textbf{one transition} (e.g.\ Fade) applied consistently to the whole presentation; use animations only when they \textbf{serve the message} (progressive reveal, chart by series); prefer fast, sober effects without sound; and always rehearse to check that the effects do not slow the delivery.

\emph{Conclusion.} Transitions and animations improve a presentation only when they are purposeful, consistent and discreet. The statement is therefore false in its absolute form: the effect must serve the content, never replace it.

\textbf{Examiner's expectations:} a genuine discussion (both sides), precise vocabulary, concrete examples, and a justified personal conclusion — a purely one-sided essay cannot earn full marks.
\end{corrige}

\newpage

\coursec{Summary sheet}

\begin{popfigure}
\begin{tikzpicture}[
  mindmap,
  concept color=popblue,
  level 1/.style={level distance=4.5cm,sibling angle=60},
  level 2/.style={level distance=3.2cm,sibling angle=45},
  every node/.style={font=\small, text width=3.2cm, align=center},
  concept/.append style={fill=popblue!20, draw=popblue, thick},
  child concept/.append style={fill=popgreen!20, draw=popgreen, thick},
  grandchild concept/.append style={fill=poporange!20, draw=poporange, thick}
]
\node[concept, align=center]{ICT22\\Producing Slides}
  child[concept color=popgreen] {node[child concept, align=center]{Presentation\\Software Basics}
    child[concept color=poporange] {node[grandchild concept, align=center]{Interface\\ \& Tools}}
    child[concept color=poporange] {node[grandchild concept, align=center]{File Formats\\ (pptx, odp)}}
  }
  child[concept color=poppurple] {node[child concept, align=center]{Slide Design\\Principles}
    child[concept color=poppink] {node[grandchild concept, align=center]{Layout \&\\Consistency}}
    child[concept color=poppink] {node[grandchild concept, align=center]{Typography\\ \& Colour}}
  }
  child[concept color=popblue] {node[child concept, align=center]{Multimedia\\Integration}
    child[concept color=popgold] {node[grandchild concept, align=center]{Images,\\ Audio, Video}}
    child[concept color=popgold] {node[grandchild concept, align=center]{Charts,\\ SmartArt}}
  }
  child[concept color=poporange] {node[child concept, align=center]{Transitions\\ \& Animations}
    child[concept color=popgreen] {node[grandchild concept, align=center]{Slide\\Transitions}}
    child[concept color=popgreen] {node[grandchild concept, align=center]{Object\\Animations}}
  }
  child[concept color=poppink] {node[child concept, align=center]{Delivery,\\Printing \& Export}
    child[concept color=poppurple] {node[grandchild concept, align=center]{Presenter\\View}}
    child[concept color=poppurple] {node[grandchild concept, align=center]{PDF/Handout\\Export}}
  }
  child[concept color=popgold] {node[child concept, align=center]{Evaluation\\ \& Best Practices}
    child[concept color=popblue] {node[grandchild concept, align=center]{Audience\\Feedback}}
    child[concept color=popblue] {node[grandchild concept, align=center]{Accessibility\\Checks}}
  };
\end{tikzpicture}
\legende{Mind map of ICT22: Producing Slides that Display Illustrative and Attractive Presentations}
\end{popfigure}

\begin{bilan}
\textbf{Key Definitions}
\begin{itemize}
  \item \cle{Presentation software}: Application (e.g. LibreOffice Impress, Microsoft PowerPoint) used to create, edit and display slide decks.
  \item \cle{Slide master}: Template that defines default layout, fonts, colours and placeholders for all slides.
  \item \cle{Transition}: Visual effect applied when moving from one slide to the next.
  \item \cle{Animation}: Effect applied to individual objects (text, images, charts) within a slide.
  \item \cle{Multimedia}: Integration of non‑text media — images, audio, video, charts, SmartArt — into slides.
  \item \cle{Presenter view}: Mode showing current slide, next slide, notes and timer on the speaker’s screen only.
\end{itemize}

\textbf{Laws / Principles / Formulas (Examinable)}
\begin{itemize}
  \item \textbf{Consistency rule}: Use a single slide master; limit fonts to two families (heading + body) and a palette of 3–4 colours.
  \item \textbf{6×6 guideline}: Maximum six bullet points per slide, six words per bullet.
  \item \textbf{Contrast principle}: Text/background contrast ratio ≥ 4.5:1 (WCAG AA) for readability.
  \item \textbf{Animation hierarchy}: Entrance → Emphasis → Exit; avoid simultaneous multiple animations on one object.
  \item \textbf{File size optimisation}: Compress images (≤ 150 dpi for screen, ≤ 300 dpi for print); embed only used fonts.
\end{itemize}

\textbf{Methods (Step‑by‑Step)}
\begin{methode}[Creating a Consistent Slide Deck]
\begin{itemize}
  \item Open the software and choose a blank presentation.
  \item Define a \cle{Slide Master}: set theme colours, heading/body fonts, logo placeholder, footer (date, slide number).
  \item Insert slides using predefined layouts (Title, Title+Content, Two Content, etc.).
  \item Apply the 6×6 guideline; keep text concise, use keywords.
  \item Insert multimedia via \texttt{Insert → Media}; resize proportionally (hold Shift).
  \item Add transitions: \texttt{Transitions tab → Choose effect → Set duration → Apply to All} (or per slide).
  \item Animate objects: \texttt{Animations tab → Add Animation → Set Start (On Click/After Previous) → Adjust timing}.
  \item Rehearse with \cle{Presenter View}; check timing, notes, laser pointer.
  \item Export: \texttt{File → Export → Create PDF/XPS} for handouts; \texttt{Package Presentation for CD} for offline delivery.
\end{itemize}
\end{methode}

\begin{methode}[Evaluating a Presentation]
\begin{itemize}
  \item Run a peer review: checklist for clarity, visual hierarchy, spelling, accessibility (alt text for images).
  \item Collect audience feedback (Google Form, Mentimeter) on content relevance and delivery pace.
  \item Record timings; adjust slide count or animation speed if over/under allotted time.
\end{itemize}
\end{methode}

\textbf{Common Mistakes (Traps)}
\begin{attention}[Frequent Errors]
\begin{itemize}
  \item Overloading slides with paragraphs instead of bullet points.
  \item Using more than three font families or clashing colour schemes.
  \item Applying different transitions on every slide → distracts audience.
  \item Animating every object; results in chaotic, unprofessional effect.
  \item Forgetting to embed fonts or compress images → huge file, playback issues on other machines.
  \item Neglecting alt text for images → fails accessibility checks.
  \item Not testing Presenter View on the actual projection hardware before the exam/presentation.
\end{itemize}
\end{attention}

\textbf{GCE Advanced Level Tips}
\begin{aretenir}[Exam‑Focused Advice]
\begin{itemize}
  \item \textbf{Practical paper}: You will be asked to produce a 6–8 slide deck in 45 min. Master the Slide Master first; it saves time.
  \item \textbf{Theory questions}: Define transition vs animation; state the 6×6 rule; explain why compressing images matters.
  \item \textbf{Short‑answer}: “Give two ways to ensure a presentation is accessible.” → Alt text, high contrast, readable font size (≥ 24 pt).
  \item \textbf{Scenario}: “A slide contains a 5 MB video. The file must be emailed (limit 10 MB). What do you do?” → Compress video, link to online storage, or convert to lower resolution.
  \item \textbf{Keyboard shortcuts}: \texttt{Ctrl+M} (new slide), \texttt{F5} (start show), \texttt{Shift+F5} (start from current), \texttt{Ctrl+P} (pen in show).
  \item \textbf{Cameroon context}: Mention using MTN/Orange mobile data to upload to Google Drive for collaboration; reference ANTIC guidelines for secure presentation sharing in e‑government.
\end{itemize}
\end{aretenir}
\end{bilan}

\findecours

\end{document}
